\begin{theorem}[离散 kink 扇的连续折线化]\label{thm:conclusion-discrete-kink-fan-piecewise-linearization}
设
\[
b:\{2,3,\dots,Q\}\to\RR
\]
为离散凸函数。定义其分段仿射插值
\[
\bar b(t):=(q+1-t)b(q)+(t-q)b(q+1),
\qquad
t\in[q,q+1],\quad 2\le q\le Q-1.
\]
则 \(\bar b\) 为凸函数。并且对每个内点 \(q\in\{3,\dots,Q-1\}\)，其凸次微分满足
\[
\partial \bar b(q)=\bigl[b(q)-b(q-1),\ b(q+1)-b(q)\bigr].
\]
特别地，\(q\) 是离散意义下的 kink，当且仅当 \(\bar b\) 在 \(q\) 处不可微。

进一步，对任意 \(\gamma\in\RR\)，定义
\[
F_\gamma(q):=(q-1)\gamma-b(q),
\qquad
\bar F_\gamma(t):=(t-1)\gamma-\bar b(t).
\]
则有严格等价
\[
q\in \operatorname*{argmax}_{2\le r\le Q}F_\gamma(r)
\iff
q\in \operatorname*{argmax}_{2\le t\le Q}\bar F_\gamma(t).
\]
\end{theorem}
\begin{proof}
离散凸性意味着相邻离散斜率
\[
b(q+1)-b(q)
\]
随 \(q\) 单调不减；而 \(\bar b\) 在每个区间 \([q,q+1]\) 上恰以此为常斜率，故 \(\bar b\) 凸。内点 \(q\) 处的左、右导数分别为
\[
\bar b'_-(q)=b(q)-b(q-1),
\qquad
\bar b'_+(q)=b(q+1)-b(q),
\]
于是
\[
\partial\bar b(q)=[\bar b'_-(q),\bar b'_+(q)].
\]
因此 \(\bar b\) 在 \(q\) 处不可微当且仅当两端斜率不等，即 \(q\) 为 kink。

再看最优化。对每个整数区间 \([q,q+1]\)，\(\bar F_\gamma\) 为仿射函数，故其在该区间上的最大值必取于端点；若斜率为零，则整个区间皆为极大集，但端点仍为极大点。于是 \(\bar F_\gamma\) 的全局极大点集合总包含某个整数点，且其整数极大点恰与 \(F_\gamma\) 的极大点相同。证毕。
\end{proof}
\leanverified{paper\_conclusion\_discrete\_kink\_fan\_piecewise\_linearization}

\begin{theorem}[连续压力 secant-majorant 的离散 Legendre 恒等式]\label{thm:conclusion-continuous-pressure-secant-majorant-discrete-legendre}
设
\[
\Lambda:[1,N]\to\RR
\]
为凸函数。以整数节点 \(1,2,\dots,N\) 构造其 secant-majorant
\[
\Lambda_N^{\mathrm{sec}}(t):=(n+1-t)\Lambda(n)+(t-n)\Lambda(n+1),
\qquad
t\in[n,n+1],\quad 1\le n\le N-1.
\]
则有
\[
\Lambda_N^{\mathrm{sec}}(t)\ge \Lambda(t)
\qquad
(1\le t\le N),
\]
并且其 Legendre--Fenchel 变换满足精确闭式
\[
\bigl(\Lambda_N^{\mathrm{sec}}\bigr)^\ast(\gamma)
:=
\sup_{1\le t\le N}\bigl(t\gamma-\Lambda_N^{\mathrm{sec}}(t)\bigr)
=
\sup_{1\le n\le N}\bigl(n\gamma-\Lambda(n)\bigr).
\]
\end{theorem}
\begin{proof}
凸函数的图像位于任意两点连线之下，故 \(\Lambda_N^{\mathrm{sec}}\ge\Lambda\)。对每个区间 \([n,n+1]\)，函数
\[
t\longmapsto t\gamma-\Lambda_N^{\mathrm{sec}}(t)
\]
为仿射函数，因此其区间上极大值必取于端点。对所有区间取并，即得
\[
\sup_{1\le t\le N}\bigl(t\gamma-\Lambda_N^{\mathrm{sec}}(t)\bigr)
=
\sup_{1\le n\le N}\bigl(n\gamma-\Lambda(n)\bigr).
\]
证毕。
\end{proof}

\begin{theorem}[离散 \(q\)-扇即连续压力折线包的 exposed fan]\label{thm:conclusion-discrete-q-fan-exposed-fan-of-pressure-secant}
\leanverified{paper\_conclusion\_discrete\_q\_fan\_exposed\_fan\_of\_pressure\_secant}
令附录 \ref{app:pom-projective-pressure-operator} 中的连续压力记为 \(\Lambda(t)\)，并在有限窗 \(q\in\{2,\dots,Q\}\) 上定义
\[
b_\Lambda(q):=\Lambda(q-1).
\]
则由整数点退化关系
\[
\Lambda(q-1)=\log r_q-\log|A_{\mathrm{out}}|
\]
可知，\(b_\Lambda\) 正是连续压力在整点 \(t=q-1\) 上的采样。对任意内点 \(q\in\{3,\dots,Q-1\}\)，有精确相带
\[
q\in \operatorname*{argmax}_{2\le r\le Q}\bigl((r-1)\gamma-b_\Lambda(r)\bigr)
\iff
\Lambda(q-1)-\Lambda(q-2)\le \gamma\le \Lambda(q)-\Lambda(q-1).
\]
并且 \(q\) 属于 kink 集，当且仅当整数节点 \(t=q-1\) 是折线包 \(\Lambda_{Q-1}^{\mathrm{sec}}\) 的一个不可微点。
\end{theorem}
\begin{proof}
由注 \ref{rem:pom-projective-pressure-route}、定理 \ref{thm:pom-projective-pressure-zero-normalization} 与定理 \ref{thm:pom-projective-operator-degenerates-to-moment-kernel}，
\[
\Lambda(q-1)=\log r_q-\log|A_{\mathrm{out}}|
\qquad(q\ge 1).
\]
故 \(b_\Lambda\) 正是连续压力在整数节点上的采样。将定理 \ref{thm:conclusion-discrete-kink-fan-piecewise-linearization} 应用于
\[
b(q)=b_\Lambda(q)=\Lambda(q-1)
\]
即可得到
\[
\partial^- b_\Lambda(q)=\Lambda(q-1)-\Lambda(q-2),
\qquad
\partial^+ b_\Lambda(q)=\Lambda(q)-\Lambda(q-1),
\]
从而得到所述 slope-band 判据。最后，\(q\) 为离散 kink 当且仅当 \(b_\Lambda\) 的分段仿射插值在 \(q\) 处不可微，这等价于 \(\Lambda_{Q-1}^{\mathrm{sec}}\) 在节点 \(t=q-1\) 处不可微。证毕。
\end{proof}

\begin{theorem}[离散 Chernoff 扇对连续 Legendre 前沿的严格夹持]\label{thm:conclusion-discrete-chernoff-fan-strictly-sandwiched-by-continuous-legendre}
在结论 \ref{thm:conclusion-discrete-q-fan-exposed-fan-of-pressure-secant} 的设定下，定义有限窗离散前沿
\[
J_Q(\gamma)
:=
\sup_{2\le q\le Q}\bigl((q-1)\gamma-\Lambda(q-1)\bigr)
=
\sup_{1\le n\le Q-1}\bigl(n\gamma-\Lambda(n)\bigr),
\]
以及相应连续前沿
\[
I_Q(\gamma):=
\sup_{1\le t\le Q-1}\bigl(t\gamma-\Lambda(t)\bigr).
\]
则有严格夹持
\[
J_Q(\gamma)
=
\bigl(\Lambda_{Q-1}^{\mathrm{sec}}\bigr)^\ast(\gamma)
\le
I_Q(\gamma).
\]
若 \(I_Q(\gamma)\) 的某个极大点落在整数节点，或更一般地落在某个单位区间上而 \(\Lambda\) 在该区间内恰为仿射，则有
\[
J_Q(\gamma)=I_Q(\gamma).
\]
反之，若 \(I_Q(\gamma)\) 的全部极大点都落在某个开区间 \((n,n+1)\) 内，且 \(\Lambda\) 在该区间上严格低于其弦线，则
\[
J_Q(\gamma)<I_Q(\gamma).
\]
\end{theorem}
\begin{proof}
第一式由定理 \ref{thm:conclusion-continuous-pressure-secant-majorant-discrete-legendre} 立即得到。又因
\[
\Lambda_{Q-1}^{\mathrm{sec}}\ge \Lambda,
\]
故
\[
t\gamma-\Lambda_{Q-1}^{\mathrm{sec}}(t)\le t\gamma-\Lambda(t)
\]
逐点成立，取上确界即得
\[
J_Q(\gamma)\le I_Q(\gamma).
\]
若某个连续极大点 \(t_\ast\) 满足
\[
\Lambda_{Q-1}^{\mathrm{sec}}(t_\ast)=\Lambda(t_\ast),
\]
则在该点两种目标函数取同值；由于左端本就不超过右端，遂得等号。整数节点与区间仿射两种情形都满足这一条件。

若反之每个连续极大点都落在某个开区间上且严格有
\[
\Lambda_{Q-1}^{\mathrm{sec}}(t_\ast)>\Lambda(t_\ast),
\]
则在每个连续极大点处左目标严格小于右目标，因而上确界亦严格小于，得 \(J_Q(\gamma)<I_Q(\gamma)\)。证毕。
\end{proof}

\begin{theorem}[整数相的精确可达性判据]\label{thm:conclusion-integer-phase-exact-attainability}
固定 \(n\in\{1,\dots,Q-1\}\)。则以下三条件等价：
\[
n\in \operatorname*{argmax}_{1\le t\le Q-1}\bigl(t\gamma-\Lambda(t)\bigr);
\]
\[
\gamma\in \partial \Lambda(n);
\]
\[
J_Q(\gamma)=I_Q(\gamma)=n\gamma-\Lambda(n).
\]
并且总有包含关系
\[
\partial \Lambda(n)\subseteq
\bigl[\Lambda(n)-\Lambda(n-1),\ \Lambda(n+1)-\Lambda(n)\bigr]
\qquad
(1<n<Q-1).
\]
\leanverified{paper\_conclusion\_integer\_phase\_exact\_attainability}
\end{theorem}
\begin{proof}
由凸分析的标准切线判据，
\[
n\in\operatorname*{argmax}_{1\le t\le Q-1}\bigl(t\gamma-\Lambda(t)\bigr)
\iff
\Lambda(t)\ge \Lambda(n)+\gamma(t-n)\quad \forall t,
\]
右式正是 \(\gamma\in\partial\Lambda(n)\) 的定义，故前两式等价。

若 \(\gamma\in\partial\Lambda(n)\)，则 \(n\) 同时是连续前沿的极大点；而 \(n\) 亦在离散取值集中，于是
\[
J_Q(\gamma)\ge n\gamma-\Lambda(n)=I_Q(\gamma).
\]
再由定理 \ref{thm:conclusion-discrete-chernoff-fan-strictly-sandwiched-by-continuous-legendre} 的一般不等式 \(J_Q(\gamma)\le I_Q(\gamma)\)，遂得三式同时成立。

最后，对凸函数 \(\Lambda\)，相邻节点上的差商夹住次微分：
\[
\Lambda(n)-\Lambda(n-1)\le g\le \Lambda(n+1)-\Lambda(n)
\qquad
(\forall g\in\partial\Lambda(n)),
\]
故得所示包含关系。证毕。
\end{proof}

\begin{theorem}[幽灵扇区]\label{thm:conclusion-ghost-sector-for-discrete-secant-fan}
对内点 \(q\in\{3,\dots,Q-1\}\)，定义其离散 secant 相带
\[
B_q:=
\bigl[\Lambda(q-1)-\Lambda(q-2),\ \Lambda(q)-\Lambda(q-1)\bigr].
\]
若
\[
\partial \Lambda(q-1)\subsetneq B_q,
\]
则差集
\[
G_q:=B_q\setminus \partial \Lambda(q-1)
\]
为非空的幽灵扇区。对任意
\[
\gamma\in \operatorname{int}(G_q),
\]
有
\[
q\in \operatorname*{argmax}_{2\le r\le Q}\bigl((r-1)\gamma-\Lambda(r-1)\bigr),
\]
\[
q-1\notin \operatorname*{argmax}_{1\le t\le Q-1}\bigl(t\gamma-\Lambda(t)\bigr),
\]
并且
\[
J_Q(\gamma)<I_Q(\gamma).
\]
换言之，离散 \(q\)-相图在该扇区内仍选中相 \(q\)，但连续 Legendre 前沿已不再承认这一相；这是一种由整点采样造成的伪暴露相。
\end{theorem}
\begin{proof}
第一式由定理 \ref{thm:conclusion-discrete-q-fan-exposed-fan-of-pressure-secant} 的 secant 相带判据立即得到。第二式由
\[
\gamma\notin\partial\Lambda(q-1)
\]
与定理 \ref{thm:conclusion-integer-phase-exact-attainability} 的等价性得到。

既然 \(q-1\) 不是连续前沿的极大点，而离散前沿却仍以该节点为极大点之一，则连续极大值不可能由整数节点实现；依定理 \ref{thm:conclusion-discrete-chernoff-fan-strictly-sandwiched-by-continuous-legendre}，只能有严格不等式
\[
J_Q(\gamma)<I_Q(\gamma).
\]
证毕。
\end{proof}

\begin{theorem}[residue 细分对离散 Legendre 前沿的平移--抬升律]\label{thm:conclusion-residue-refinement-discrete-legendre-shift-lift-law}
沿用定理 \ref{thm:conclusion-residue-refinement-purity-uniformization-endpoint-rigidity} 的记号。固定有限 residue 轴大小 \(P\)，并在有限尺度 \(m\) 上定义
\[
b_m(q):=\frac1m\log S_q(m),
\qquad
b_m^{(P)}(q):=\frac1m\log S_q^{(P)}(m),
\qquad q\ge2.
\]
则对任意有限窗 \(q\in\{2,\dots,Q\}\) 与任意 \(\gamma\in\RR\)，离散前沿
\[
J_{m,Q}(\gamma):=\sup_{2\le q\le Q}\bigl((q-1)\gamma-b_m(q)\bigr),
\]
\[
J_{m,Q}^{(P)}(\gamma):=\sup_{2\le q\le Q}\bigl((q-1)\gamma-b_m^{(P)}(q)\bigr)
\]
满足双边夹持
\[
J_{m,Q}(\gamma)\le J_{m,Q}^{(P)}(\gamma)
\le
J_{m,Q}\!\left(\gamma+\frac{\log P}{m}\right).
\]
特别地：
\[
J_{m,Q}^{(P)}(\gamma)=J_{m,Q}(\gamma)\quad(\forall \gamma)
\]
当且仅当每条纤维在 residue 轴上纯；而
\[
J_{m,Q}^{(P)}(\gamma)=
J_{m,Q}\!\left(\gamma+\frac{\log P}{m}\right)
\quad(\forall \gamma)
\]
当且仅当每条纤维在 residue 轴上完全均匀化。
\end{theorem}
\begin{proof}
由定理 \ref{thm:conclusion-residue-refinement-purity-uniformization-endpoint-rigidity}，
\[
P^{1-q}S_q(m)\le S_q^{(P)}(m)\le S_q(m),
\]
故
\[
b_m(q)-\frac{q-1}{m}\log P
\le
b_m^{(P)}(q)
\le
b_m(q).
\]
由右侧不等式，
\[
(q-1)\gamma-b_m^{(P)}(q)\ge (q-1)\gamma-b_m(q),
\]
对 \(q\) 取上确界即得
\[
J_{m,Q}^{(P)}(\gamma)\ge J_{m,Q}(\gamma).
\]
由左侧不等式，
\[
(q-1)\gamma-b_m^{(P)}(q)
\le
(q-1)\left(\gamma+\frac{\log P}{m}\right)-b_m(q),
\]
再取上确界得
\[
J_{m,Q}^{(P)}(\gamma)
\le
J_{m,Q}\!\left(\gamma+\frac{\log P}{m}\right).
\]

若每条纤维在 residue 轴上纯，则
\[
S_q^{(P)}(m)=S_q(m)
\qquad(\forall q>1),
\]
故两前沿完全相同。若每条纤维在 residue 轴上完全均匀化，则
\[
S_q^{(P)}(m)=P^{1-q}S_q(m),
\]
从而
\[
b_m^{(P)}(q)=b_m(q)-\frac{q-1}{m}\log P,
\]
代回定义即得精确平移公式。证毕。
\end{proof}

\begin{theorem}[证书前沿恰为离散率函数的广义逆表述]\label{thm:conclusion-discrete-certificate-front-generalized-inverse}
\leanverified{paper\_conclusion\_discrete\_certificate\_front\_generalized\_inverse}
沿用附录 \ref{app:pom-projective-pressure-operator} 的记号，并在已核验指数口径
\[
S_q(m)=\Theta(r_q^m)
\]
下记
\[
\Phi(q):=\tau(q)-\log 2=\log r_q-\log 2,
\qquad
c_m(\varepsilon):=\frac1m\log\frac1\varepsilon
\qquad
(q\ge 2).
\]
定义离散率函数
\[
I_{\mathrm{disc}}(\gamma):=\sup_{q\ge 2}\bigl((q-1)\gamma-\Phi(q)\bigr),
\]
以及证书前沿
\[
\gamma_m^{\mathrm{cert}}(\varepsilon):=
\inf_{q\ge 2}\frac{\Phi(q)+c_m(\varepsilon)}{q-1}.
\]
则对任意 \(m\ge 1\) 与 \(\varepsilon\in(0,1)\)，有严格恒等式
\[
\gamma_m^{\mathrm{cert}}(\varepsilon)
=
\inf\bigl\{\gamma\ge 0:\ I_{\mathrm{disc}}(\gamma)\ge c_m(\varepsilon)\bigr\}.
\]
换言之，prime-register 的最省指数预算正是离散 Legendre--Fenchel 率函数在高度 \(c_m(\varepsilon)\) 处的最左水平截线。
\end{theorem}
\begin{proof}
由定义，
\[
I_{\mathrm{disc}}(\gamma)\ge c_m(\varepsilon)
\iff
\sup_{q\ge 2}\bigl((q-1)\gamma-\Phi(q)\bigr)\ge c_m(\varepsilon).
\]
这又等价于存在某个 \(q\ge 2\) 使
\[
(q-1)\gamma-\Phi(q)\ge c_m(\varepsilon),
\]
亦即
\[
\gamma\ge \frac{\Phi(q)+c_m(\varepsilon)}{q-1}.
\]
因此满足 \(I_{\mathrm{disc}}(\gamma)\ge c_m(\varepsilon)\) 的全部 \(\gamma\) 恰构成半直线
\[
\Bigl[
\inf_{q\ge2}\frac{\Phi(q)+c_m(\varepsilon)}{q-1},
\ \infty
\Bigr),
\]
其左端点正是 \(\gamma_m^{\mathrm{cert}}(\varepsilon)\)。证毕。
\end{proof}

\begin{corollary}[证书前沿关于置信代价的凹折线结构]\label{cor:conclusion-certificate-front-concave-polygonal-kink-support}
令 \(c\ge 0\) 为抽象置信代价，并定义
\[
\Gamma_q(c):=\frac{\Phi(q)+c}{q-1},
\qquad
G(c):=\inf_{q\ge 2}\Gamma_q(c).
\]
则 \(G\) 是非减、分段仿射、凹函数，并且每一段的斜率必属于集合
\[
\left\{\frac1{q-1}:q\ge 2\right\}.
\]
若某个 \(q\) 在 \(G\) 的一段开区间上处于活跃地位，则 \(q\) 必属于离散凸预算的 kink 集。并且对任意 \(q<r\)，两条候选分支的交点显式为
\[
c_{q,r}=
\frac{(q-1)\Phi(r)-(r-1)\Phi(q)}{r-q}.
\]
因此，证书前沿随置信代价 \(c\) 的全部换相由 kink 分支拼接而成，而不再是对整数矩阶的离散穷举。
\end{corollary}
\begin{proof}
每个 \(\Gamma_q(c)\) 都是关于 \(c\) 的仿射函数，斜率为 \(1/(q-1)\)。故其下包络 \(G(c)\) 必为凹的分段仿射函数，且斜率只能取自上述集合。

若某个 \(q\) 在某个开区间 \(J\subset[0,\infty)\) 上始终活跃，则对每个 \(c\in J\) 有
\[
G(c)=\Gamma_q(c)=\frac{\Phi(q)+c}{q-1}.
\]
记 \(\gamma:=G(c)\)。则
\[
c=(q-1)\gamma-\Phi(q).
\]
又因为 \(\Gamma_q(c)\le \Gamma_r(c)\) 对一切 \(r\ge 2\) 成立，故
\[
\gamma\le \frac{\Phi(r)+c}{r-1}
\iff
(r-1)\gamma-\Phi(r)\le c=(q-1)\gamma-\Phi(q).
\]
于是 \(q\) 在 \(\gamma\) 处实现 \(I_{\mathrm{disc}}(\gamma)\) 的上确界。由于 \(J\) 为开区间，而 \(\gamma=\Gamma_q(c)\) 在 \(J\) 上亦跑过一个开区间，故 \(q\) 在一个开 \(\gamma\)-区间上持续支撑离散率函数。依定理 \ref{thm:conclusion-discrete-convex-budget-kink-fan}，只有 kink 才能在开区间上稳定支撑；非 kink 只能在退化单点上并列出现。

最后，\(\Gamma_q(c)=\Gamma_r(c)\) 直接解出
\[
\frac{\Phi(q)+c}{q-1}=\frac{\Phi(r)+c}{r-1}
\iff
c=\frac{(q-1)\Phi(r)-(r-1)\Phi(q)}{r-q}.
\]
证毕。
\end{proof}
\leanverified{paper\_conclusion\_certificate\_front\_concave\_polygonal\_kink\_support}

\begin{theorem}[连续压力前沿对整数矩阶证书的严格支配]\label{thm:conclusion-continuous-pressure-front-dominates-discrete-certificate}
设附录 \ref{app:pom-projective-pressure-operator} 的连续压力曲线
\[
\Lambda:[0,\infty)\to\RR
\]
已存在，并满足整点一致性
\[
\Lambda(q-1)=\Phi(q)
\qquad
(q\ge 2).
\]
定义连续率函数
\[
I(\gamma):=\sup_{t\ge 0}\bigl(t\gamma-\Lambda(t)\bigr),
\]
以及连续证书前沿
\[
\gamma_m^{\mathrm{cont}}(\varepsilon):=
\inf_{t>0}\frac{\Lambda(t)+c_m(\varepsilon)}{t}.
\]
则有
\[
\gamma_m^{\mathrm{cont}}(\varepsilon)
=
\inf\bigl\{\gamma\ge 0:\ I(\gamma)\ge c_m(\varepsilon)\bigr\},
\]
并且
\[
I(\gamma)\ge I_{\mathrm{disc}}(\gamma)
\qquad(\forall \gamma\ge 0),
\qquad
\gamma_m^{\mathrm{cont}}(\varepsilon)\le \gamma_m^{\mathrm{cert}}(\varepsilon).
\]
特别地，连续压力前沿从不劣于任意整数矩阶证书；它是离散 Chernoff 家族的真正凸化闭包。
\leanverified{paper\_conclusion\_continuous\_pressure\_front\_dominates\_discrete\_certificate}
\end{theorem}
\begin{proof}
第一式与定理 \ref{thm:conclusion-discrete-certificate-front-generalized-inverse} 的证明完全同型，只需把离散指标 \(q-1\) 替换为连续变量 \(t>0\)。

又由于整点集合 \(\{q-1:q\ge 2\}\subset[0,\infty)\)，有
\[
I(\gamma)=\sup_{t\ge 0}(t\gamma-\Lambda(t))
\ge
\sup_{q\ge 2}\bigl((q-1)\gamma-\Lambda(q-1)\bigr)
=
\sup_{q\ge 2}\bigl((q-1)\gamma-\Phi(q)\bigr)
=
I_{\mathrm{disc}}(\gamma).
\]
对两边的广义逆取最左端点，即得
\[
\gamma_m^{\mathrm{cont}}(\varepsilon)\le \gamma_m^{\mathrm{cert}}(\varepsilon).
\]
证毕。
\end{proof}

\begin{theorem}[连续 Legendre 前沿的 Euler--Fenchel 接触方程]\label{thm:conclusion-continuous-legendre-euler-fenchel-contact}
进一步设 \(\Lambda\) 在 \((0,\infty)\) 上可微且本质严格凸。若
\[
t_{m,\varepsilon}\in(0,\infty)
\]
是函数
\[
F_{m,\varepsilon}(t):=\frac{\Lambda(t)+c_m(\varepsilon)}{t}
\]
的内点极小值，则必满足精确接触方程
\[
t_{m,\varepsilon}\Lambda'(t_{m,\varepsilon})-\Lambda(t_{m,\varepsilon})
=
c_m(\varepsilon).
\]
若记
\[
\gamma_{m,\varepsilon}^{\ast}:=\Lambda'(t_{m,\varepsilon}),
\]
则
\[
I(\gamma_{m,\varepsilon}^{\ast})=c_m(\varepsilon),
\qquad
\gamma_{m,\varepsilon}^{\ast}
=
\frac{\Lambda(t_{m,\varepsilon})+c_m(\varepsilon)}{t_{m,\varepsilon}}
=
\gamma_m^{\mathrm{cont}}(\varepsilon).
\]
换言之，连续 prime-register 前沿由一条 Legendre 曲线与水平代价 \(c_m(\varepsilon)\) 的切触方程唯一生成。
\end{theorem}
\begin{proof}
对 \(F_{m,\varepsilon}\) 求导，
\[
F_{m,\varepsilon}'(t)
=
\frac{t\Lambda'(t)-\Lambda(t)-c_m(\varepsilon)}{t^2}.
\]
在内点极小值 \(t=t_{m,\varepsilon}\) 处，必有 \(F_{m,\varepsilon}'(t_{m,\varepsilon})=0\)，于是
\[
t_{m,\varepsilon}\Lambda'(t_{m,\varepsilon})-\Lambda(t_{m,\varepsilon})
=
c_m(\varepsilon).
\]
置 \(\gamma_{m,\varepsilon}^{\ast}:=\Lambda'(t_{m,\varepsilon})\)。由于 \(\Lambda\) 本质严格凸，\(t_{m,\varepsilon}\) 正是 \(I(\gamma_{m,\varepsilon}^{\ast})\) 的唯一对偶接触点，因此
\[
I(\gamma_{m,\varepsilon}^{\ast})
=
t_{m,\varepsilon}\gamma_{m,\varepsilon}^{\ast}-\Lambda(t_{m,\varepsilon})
=
c_m(\varepsilon).
\]
再由上式整理得
\[
\gamma_{m,\varepsilon}^{\ast}
=
\frac{\Lambda(t_{m,\varepsilon})+c_m(\varepsilon)}{t_{m,\varepsilon}}
=
F_{m,\varepsilon}(t_{m,\varepsilon})
=
\gamma_m^{\mathrm{cont}}(\varepsilon).
\]
证毕。
\leanverified{paper\_conclusion\_continuous\_legendre\_euler\_fenchel\_contact}
\end{proof}

\begin{corollary}[固定置信度下连续证书指数向典型纤维指数收敛]\label{cor:conclusion-fixed-confidence-continuous-certificate-typical-fiber-limit}
设 \(I\) 为良函数，且其唯一零点为
\[
\gamma_{\ast}=\Lambda'(0).
\]
对任意固定 \(\varepsilon\in(0,1)\)，有
\[
\gamma_m^{\mathrm{cont}}(\varepsilon)\downarrow \gamma_{\ast}
\qquad (m\to\infty).
\]
若定义连续前沿所对应的最小寄存器预算为
\[
P_m^{\mathrm{cont}}(\varepsilon):=
2\exp\bigl(m\gamma_m^{\mathrm{cont}}(\varepsilon)\bigr),
\]
则
\[
\lim_{m\to\infty}\frac1m\log P_m^{\mathrm{cont}}(\varepsilon)=\gamma_{\ast}.
\]
因此，在固定失败概率口径下，最优空间预算的指数斜率最终由典型纤维指数而非最坏纤维指数支配。
\end{corollary}
\begin{proof}
因
\[
c_m(\varepsilon)=\frac1m\log\frac1\varepsilon\downarrow 0,
\]
而 \(I\) 为良凸函数，且唯一零点为 \(\gamma_{\ast}\)，故在 \([\gamma_{\ast},\infty)\) 上必有
\[
\{\gamma:I(\gamma)\ge c\}=[\Gamma(c),\infty)
\]
对某个单调函数 \(\Gamma(c)\) 成立，且 \(\Gamma(c)\downarrow \gamma_{\ast}\) 当 \(c\downarrow 0\)。由定理 \ref{thm:conclusion-continuous-pressure-front-dominates-discrete-certificate}，
\[
\gamma_m^{\mathrm{cont}}(\varepsilon)=\Gamma(c_m(\varepsilon)),
\]
故第一式成立。第二式由
\[
\frac1m\log P_m^{\mathrm{cont}}(\varepsilon)
=
\frac{\log 2}{m}+\gamma_m^{\mathrm{cont}}(\varepsilon)
\]
及 \(\log 2/m\to 0\) 立即得到。证毕。
\end{proof}

\begin{theorem}[固定置信度下连续证书预算的平方根修正律]\label{thm:conclusion-fixed-confidence-continuous-certificate-sqrtm-correction}
在结论 \ref{cor:conclusion-fixed-confidence-continuous-certificate-typical-fiber-limit} 的设定下，进一步假设 \(\Lambda\) 在 \(0\) 的邻域二次可微，并记
\[
\sigma^2:=\Lambda''(0)>0.
\]
则当 \(m\to\infty\) 时，
\[
\gamma_m^{\mathrm{cont}}(\varepsilon)
=
\gamma_{\ast}
+
\sqrt{\frac{2\sigma^2\log(1/\varepsilon)}{m}}
+
o(m^{-1/2}),
\]
从而
\[
\log P_m^{\mathrm{cont}}(\varepsilon)
=
\log 2 + m\gamma_{\ast}
+
\sqrt{2\sigma^2\,m\log(1/\varepsilon)}
+
o(\sqrt m).
\]
这表明固定置信度下的最优 prime-register 预算除线性主项 \(m\gamma_{\ast}\) 外，还具有一个由中心极限定理方差 \(\sigma^2\) 精确控制的 \(\sqrt m\)-级修正项。
\end{theorem}
\begin{proof}
由 \(\Lambda\) 在 \(0\) 的二次展开，
\[
\Lambda(t)=\gamma_{\ast}t+\frac{\sigma^2}{2}t^2+o(t^2)
\qquad (t\to 0).
\]
于是
\[
\Lambda'(t)=\gamma_{\ast}+\sigma^2 t+o(t),
\]
并且
\[
I(\gamma)=\frac{(\gamma-\gamma_{\ast})^2}{2\sigma^2}
+o\!\left((\gamma-\gamma_{\ast})^2\right)
\qquad (\gamma\downarrow\gamma_{\ast}).
\]
令 \(\gamma=\gamma_m^{\mathrm{cont}}(\varepsilon)\)。由定理 \ref{thm:conclusion-continuous-pressure-front-dominates-discrete-certificate}，
\[
I(\gamma_m^{\mathrm{cont}}(\varepsilon))=c_m(\varepsilon)=\frac1m\log\frac1\varepsilon.
\]
代入上式并反解，得到
\[
\gamma_m^{\mathrm{cont}}(\varepsilon)-\gamma_{\ast}
=
\sqrt{\frac{2\sigma^2\log(1/\varepsilon)}{m}}
+
o(m^{-1/2}).
\]
再乘以 \(m\) 并加上 \(\log 2\) 即得预算展开。证毕。
\end{proof}
\leanverified{paper\_conclusion\_fixed\_confidence\_continuous\_certificate\_sqrtm\_correction}

\begin{theorem}[有限矩阶相图的单证书编译]\label{thm:conclusion-finite-moment-order-phase-diagram-single-certificate-compilation}
固定有限窗 \(Q\ge 3\)。沿用附录 \ref{app:pom-projective-pressure-operator} 的记号，记
\[
\beta_q:=\log r_q-\log |A_{\mathrm{out}}|=\Lambda(q-1),
\qquad
q=2,\dots,Q.
\]
对任意非空指标集 \(J\subset\{2,\dots,Q\}\)，定义块对角 Perron 证书对象
\[
\mathfrak A_J:=\bigoplus_{q\in J}\widetilde A_q,
\qquad
\mathfrak u_J:=\bigoplus_{q\in J}\widetilde u_q,
\qquad
\mathfrak v_J:=\bigoplus_{q\in J}\widetilde v_q.
\]
则 \(\mathfrak A_J\) 连同其块标记唯一编码了全部整数节点
\[
\{r_q:\ q\in J\},
\]
并满足精确维数上界
\[
\dim \mathfrak A_J
=
\sum_{q\in J}\dim \widetilde A_q
\le
\sum_{q\in J}\binom{q+9}{9}.
\]
特别地，若取 \(J=\{2,\dots,Q\}\)，则整段离散压力节点 \(\{\beta_q\}_{q=2}^{Q}\) 可由一个单一块对角非负证书对象完全承载，其总维数满足
\[
\dim \mathfrak A_{\{2,\dots,Q\}}
\le
\sum_{q=2}^{Q}\binom{q+9}{9}
=
\binom{Q+10}{10}-11.
\]
\end{theorem}
\begin{proof}
由命题 \ref{prop:pom-moment-symmetric-power}，对每个整数 \(q\ge 2\) 已构造出对称幂 moment-kernel \(\widetilde A_q\) 及读出向量 \(\widetilde u_q,\widetilde v_q\)，使
\[
S_q(m)=\widetilde u_q^{\mathsf T}\widetilde A_q^{\,m-m_0}\widetilde v_q,
\qquad
\dim \widetilde A_q\le \binom{q+|Q|-1}{|Q|-1}.
\]
在主核同步状态数 \(|Q|=10\) 的口径下，
\[
\dim \widetilde A_q\le \binom{q+9}{9}.
\]
另一方面，由定理 \ref{thm:pom-projective-operator-degenerates-to-moment-kernel}，
\[
\rho(\widetilde A_q)=r_q.
\]
故对任意有限 \(J\)，取直和后得到的块对角对象 \(\mathfrak A_J\) 的第 \(q\)-块 Perron 根即为 \(r_q\)，因而 \(\mathfrak A_J\) 与块标记共同完整编码 \(\{r_q\}_{q\in J}\)。维数上界由直和可加性立即得到。

最后利用 Hockey-stick 恒等式
\[
\sum_{q=0}^{Q}\binom{q+9}{9}=\binom{Q+10}{10},
\]
减去 \(q=0,1\) 两项 \(1\) 与 \(10\)，便得
\[
\sum_{q=2}^{Q}\binom{q+9}{9}=\binom{Q+10}{10}-11.
\]
证毕。
\end{proof}
\leanverified{paper\_conclusion\_finite\_moment\_order\_phase\_diagram\_single\_certificate\_compilation}

\begin{corollary}[kink-稀疏相图的精确压缩律]\label{cor:conclusion-kink-sparse-phase-diagram-exact-compression}
沿用定理 \ref{thm:conclusion-finite-moment-order-phase-diagram-single-certificate-compilation} 的记号，并令
\[
K_Q:=\{q\in\{3,\dots,Q-1\}:\ \Delta_q^-<\Delta_q^+\}
\]
为结论 \ref{thm:conclusion-discrete-convex-budget-kink-fan} 的 stable kink 集。则整个离散相图
\[
\gamma\longmapsto \operatorname*{argmax}_{2\le q\le Q}\bigl((q-1)\gamma-\beta_q\bigr)
\]
完全由压缩证书
\[
\mathfrak A_{K_Q}=\bigoplus_{q\in K_Q}\widetilde A_q
\]
决定；所有 \(q\notin K_Q\) 的 Perron 节点在稳定相图意义下都可删去。因而离散相图的精确证书维数满足
\[
D_{\mathrm{kink}}(Q)
\le
\sum_{q\in K_Q}\binom{q+9}{9}.
\]
若记 \(\kappa_Q:=|K_Q|\)，则进一步有
\[
D_{\mathrm{kink}}(Q)\le \kappa_Q\binom{Q+9}{9}.
\]
从而
\[
\kappa_Q=o(Q)\Longrightarrow D_{\mathrm{kink}}(Q)=o(Q^{10}),
\qquad
\sup_Q\kappa_Q<\infty \Longrightarrow D_{\mathrm{kink}}(Q)=O(Q^9).
\]
\end{corollary}
\begin{proof}
结论 \ref{thm:conclusion-discrete-convex-budget-kink-fan} 已给出：离散凸预算的全部 \(q\)-选择相图本身由 kink 集生成；非 kink 阶仅可能在退化边界上并列出现，而不会形成稳定单相。将该结论应用于
\[
b(q)=\beta_q
\]
即可得到稳定相图只依赖于 \(\{\beta_q\}_{q\in K_Q}\)。而由定理 \ref{thm:conclusion-finite-moment-order-phase-diagram-single-certificate-compilation}，每个 \(\beta_q\) 又由相应的对称化 Perron 核 \(\widetilde A_q\) 完整承载。因此删去全部 \(q\notin K_Q\) 的块，不会改变相图的稳定扇分，从而得到第一条维数上界。

再由单调性
\[
\binom{q+9}{9}\le \binom{Q+9}{9}
\qquad (q\le Q)
\]
可得
\[
D_{\mathrm{kink}}(Q)\le \sum_{q\in K_Q}\binom{Q+9}{9}
=
\kappa_Q\binom{Q+9}{9}.
\]
最后利用 \(\binom{Q+9}{9}\sim Q^9/9!\) 即得两条渐近结论。证毕。
\end{proof}

\begin{theorem}[连续压力前沿的无伪整数相]\label{thm:conclusion-continuous-pressure-front-has-no-fake-integer-phases}
\leanverified{paper\_conclusion\_continuous\_pressure\_front\_has\_no\_fake\_integer\_phases}
在结论 \ref{thm:conclusion-discrete-q-fan-exposed-fan-of-pressure-secant} 的设定下，设某个内部整数节点 \(q\in\{3,\dots,Q-1\}\) 具有非空开法锥，即存在开区间 \(J\subset\RR\) 使得对每个 \(\gamma\in J\),
\[
q-1\in\operatorname*{argmax}_{t\in[1,Q-1]}\bigl(t\gamma-\Lambda(t)\bigr).
\]
则必有
\[
q\in K_Q.
\]
等价地，连续 Legendre 前沿不可能制造离散相图中并不存在的稳定整数相。
\end{theorem}
\begin{proof}
设存在 \(q\notin K_Q\) 与开区间 \(J\) 满足题设。则对任意 \(\gamma\in J\) 与任意 \(p\neq q\)，有
\[
(q-1)\gamma-\Lambda(q-1)>(p-1)\gamma-\Lambda(p-1).
\]
代入整数锚点恒等式 \(\Lambda(s-1)=\beta_s\)，得
\[
(q-1)\gamma-\beta_q>(p-1)\gamma-\beta_p
\qquad (\forall p\neq q,\ \forall \gamma\in J).
\]
这表明 \(q\) 在离散优化问题
\[
\operatorname*{argmax}_{2\le p\le Q}\bigl((p-1)\gamma-\beta_p\bigr)
\]
中是一个开 \(\gamma\)-区间上的唯一极大点。按结论 \ref{thm:conclusion-discrete-convex-budget-kink-fan}，这只可能发生在 \(q\) 为 kink 的情形，矛盾。故 \(q\in K_Q\)。证毕。
\end{proof}

\begin{theorem}[整数锚点的局部凸插值夹逼]\label{thm:conclusion-integer-anchor-local-convex-interpolation-squeeze}
\leanverified{paper\_conclusion\_integer\_anchor\_local\_convex\_interpolation\_squeeze}
在结论 \ref{thm:conclusion-continuous-pressure-front-has-no-fake-integer-phases} 的设定下，对每个 \(q=3,\dots,Q-1\) 与每个
\[
t\in[q-1,q],
\]
成立精确局部夹逼
\[
\beta_q+(t-q+1)\Delta_q^-
\le
\Lambda(t)
\le
(q-t)\beta_q+(t-q+1)\beta_{q+1}.
\]
若再有 \(q\le Q-2\)，则还可从右侧得到第二个下界
\[
\beta_{q+1}-(q-t)\Delta_{q+1}^+
\le
\Lambda(t).
\]
因此，在每个单位胞腔 \([q-1,q]\) 中，连续压力曲线被两条离散切线与一条离散弦严格夹住。
\end{theorem}
\begin{proof}
先证上界。由 \(\Lambda\) 凸且
\[
\Lambda(q-1)=\beta_q,\qquad \Lambda(q)=\beta_{q+1},
\]
对任意 \(t\in[q-1,q]\) 应有
\[
\Lambda(t)\le (q-t)\beta_q+(t-q+1)\beta_{q+1}.
\]

再证左下界。取 \(x_0=q-1\)。由于 \(\Lambda\) 凸，任一 \(g\in\partial\Lambda(x_0)\) 都给出支持线
\[
\Lambda(t)\ge \Lambda(x_0)+g(t-x_0)\qquad (\forall t).
\]
而由结论 \ref{thm:conclusion-integer-phase-exact-attainability} 的末式，
\[
\Delta_q^-=\beta_q-\beta_{q-1}\le g.
\]
对 \(t\in[q-1,q]\) 有 \(t-x_0\ge 0\)，故把 \(g\) 换成更小的 \(\Delta_q^-\) 仍保持不等式：
\[
\Lambda(t)\ge \beta_q+(t-q+1)\Delta_q^-.
\]

右侧下界同理：在 \(x_1=q\) 取某个次梯度 \(h\in\partial\Lambda(x_1)\)，有
\[
h\le \Delta_{q+1}^+,
\]
而当 \(t\le q\) 时 \(t-x_1\le 0\)，于是
\[
\Lambda(t)\ge \beta_{q+1}+h(t-q)\ge \beta_{q+1}-(q-t)\Delta_{q+1}^+.
\]
证毕。
\end{proof}

\begin{theorem}[整数锚点折线包络的法扇恰为离散 kink 扇]\label{thm:conclusion-piecewise-linear-envelope-normal-fan-equals-discrete-kink-fan}
定义整数锚点的分段线性上包络
\[
\overline\Lambda_Q:[1,Q-1]\to\RR
\]
为通过点列
\[
(q-1,\beta_q),\qquad q=2,\dots,Q
\]
的折线插值；即对 \(t\in[q-1,q]\),
\[
\overline\Lambda_Q(t):=(q-t)\beta_q+(t-q+1)\beta_{q+1}.
\]
则 \(\overline\Lambda_Q\) 是凸的，且对每个内部顶点 \(t_q:=q-1\)（\(3\le q\le Q-1\)），其法锥精确等于
\[
N_{\overline\Lambda_Q}(t_q)
=
[\Delta_q^-,\Delta_q^+].
\]
因此，\(\overline\Lambda_Q\) 的法扇分解与结论 \ref{thm:conclusion-discrete-convex-budget-kink-fan} 的离散 kink 扇分相完全一致。
\end{theorem}
\begin{proof}
由于 \(\beta_q=\Lambda(q-1)\) 来自凸函数 \(\Lambda\) 的整数采样，离散斜率单调：
\[
\Delta_q^-\le \Delta_q^+.
\]
故折线插值 \(\overline\Lambda_Q\) 本身凸。

在顶点 \(t_q=q-1\) 处，其左段斜率恰为
\[
\frac{\beta_q-\beta_{q-1}}{1}=\Delta_q^-,
\]
右段斜率恰为
\[
\frac{\beta_{q+1}-\beta_q}{1}=\Delta_q^+.
\]
对一维凸分段线性函数，顶点的法锥正是左右斜率所张成的闭区间，故
\[
N_{\overline\Lambda_Q}(t_q)=[\Delta_q^-,\Delta_q^+].
\]
而结论 \ref{thm:conclusion-discrete-convex-budget-kink-fan} 已证明：离散优化问题
\[
\operatorname*{argmax}_{2\le r\le Q}\bigl((r-1)\gamma-\beta_r\bigr)
\]
中，\(q\) 成为最优阶的充要条件正是
\[
\Delta_q^-\le \gamma\le \Delta_q^+.
\]
于是法锥区间与离散 kink 相带一一对应，结论成立。证毕。
\end{proof}

\begin{theorem}[有限节点 Legendre 审计带]\label{thm:conclusion-finite-node-legendre-audit-band}
在区间
\[
J_Q:=[2,Q-1]
\]
上，定义局部下包络
\[
\underline\Lambda_Q(t):=
\beta_q+(t-q+1)\Delta_q^-,
\qquad
t\in[q-1,q],\quad q=3,\dots,Q-1,
\]
以及上包络 \(\overline\Lambda_Q\) 如定理 \ref{thm:conclusion-piecewise-linear-envelope-normal-fan-equals-discrete-kink-fan}。则由定理 \ref{thm:conclusion-integer-anchor-local-convex-interpolation-squeeze}，
\[
\underline\Lambda_Q(t)\le \Lambda(t)\le \overline\Lambda_Q(t)
\qquad (t\in J_Q).
\]
于是对截断 Legendre 变换
\[
I_Q(\gamma):=\sup_{t\in J_Q}\bigl(t\gamma-\Lambda(t)\bigr)
\]
有严格审计带
\[
\overline I_Q(\gamma)\le I_Q(\gamma)\le \underline I_Q(\gamma),
\]
其中
\[
\overline I_Q(\gamma):=\sup_{t\in J_Q}\bigl(t\gamma-\overline\Lambda_Q(t)\bigr),
\qquad
\underline I_Q(\gamma):=\sup_{t\in J_Q}\bigl(t\gamma-\underline\Lambda_Q(t)\bigr).
\]
换言之，有限个整数 moment-kernel 节点已足以在连续大偏差前沿外侧打出一条可计算、可审计且严格的 Legendre 夹带。
\end{theorem}
\begin{proof}
由定理 \ref{thm:conclusion-integer-anchor-local-convex-interpolation-squeeze}，
\[
\underline\Lambda_Q(t)\le \Lambda(t)\le \overline\Lambda_Q(t)
\qquad (\forall t\in J_Q).
\]
对任意 \(\gamma\in\RR\)，逐点不等式乘以 \(-1\) 再加 \(t\gamma\)，得到
\[
t\gamma-\overline\Lambda_Q(t)
\le
t\gamma-\Lambda(t)
\le
t\gamma-\underline\Lambda_Q(t)
\qquad (\forall t\in J_Q).
\]
对 \(t\in J_Q\) 取上确界，便得到
\[
\overline I_Q(\gamma)\le I_Q(\gamma)\le \underline I_Q(\gamma).
\]
证毕。
\end{proof}

\begin{theorem}[离散 kink 质量的帽核曲率量子化]\label{thm:conclusion-discrete-kink-mass-hat-kernel-curvature}
\leanverified{paper\_conclusion\_discrete\_kink\_mass\_hat\_kernel\_curvature}
设 \(\tau\) 是某个开区间上的 \(C^2\) 实函数。对任意整数 \(q\)，定义离散斜率与 kink 质量
\[
\Delta\tau(q):=\tau(q)-\tau(q-1),
\qquad
\mathfrak k_\tau(q):=\Delta\tau(q+1)-\Delta\tau(q).
\]
则有精确恒等式
\[
\mathfrak k_\tau(q)
=
\int_{q-1}^{q+1}(1-|s-q|)\,\tau''(s)\,ds.
\]
特别地，若 \(\tau\) 凸，则
\[
\mathfrak k_\tau(q)=0
\iff
\tau\ \text{在 }[q-1,q+1]\text{ 上仿射}.
\]
因此，离散 kink 质量并非离散伪量，而是连续压力曲率在双胞元 \([q-1,q+1]\) 上的帽核平均。
\end{theorem}
\begin{proof}
由 Newton--Leibniz 公式，
\[
\Delta\tau(q)=\int_{q-1}^{q}\tau'(s)\,ds,
\qquad
\Delta\tau(q+1)=\int_{q}^{q+1}\tau'(s)\,ds.
\]
故
\[
\mathfrak k_\tau(q)
=
\int_0^1\bigl(\tau'(q+t)-\tau'(q-1+t)\bigr)\,dt.
\]
再把被积差写成
\[
\tau'(q+t)-\tau'(q-1+t)=\int_{q-1+t}^{q+t}\tau''(s)\,ds,
\]
并交换积分次序，得到
\[
\mathfrak k_\tau(q)
=
\int_{q-1}^{q}(s-q+1)\tau''(s)\,ds
+
\int_{q}^{q+1}(q+1-s)\tau''(s)\,ds,
\]
这正是所述帽核表示。

若 \(\tau\) 凸，则 \(\tau''\ge 0\)。当 \(\mathfrak k_\tau(q)=0\) 时，非负函数 \(\tau''\) 在严格正帽核下积分为零，从而 \(\tau''\equiv 0\) 于 \((q-1,q+1)\) 上；由连续性得 \(\tau''\equiv 0\) 于闭区间 \([q-1,q+1]\) 上，故 \(\tau\) 仿射。反向 implication 显然。证毕。
\end{proof}

\begin{theorem}[严格凸压力的满-kink 性]\label{thm:conclusion-strictly-convex-pressure-full-kink}
设 \(\tau\) 在包含整数窗 \(\{2,\dots,Q_{\max}\}\) 的开区间上可微且严格凸，并令
\[
b(q):=\tau(q),
\qquad
q\in\{2,\dots,Q_{\max}\}.
\]
则每个内部整数点
\[
q\in\{3,\dots,Q_{\max}-1\}
\]
都是 \(b\) 的 kink；若按定理 \ref{thm:conclusion-discrete-convex-budget-kink-fan} 的单边约定把端点也视为稳定相，则
\[
\operatorname{Kink}(b)=\{2,3,\dots,Q_{\max}\}.
\]
换言之，在真正的解析严格凸制度下，离散优化的 kink 稀疏化只可能发生在局部仿射区或相共存区，而不可能发生在曲率真实存在的单相区。
\end{theorem}
\begin{proof}
对任意内部整数 \(q\)，有
\[
\Delta b(q+1)-\Delta b(q)
=
\int_{q}^{q+1}\tau'(s)\,ds-\int_{q-1}^{q}\tau'(s)\,ds
=
\int_0^1\bigl(\tau'(q+t)-\tau'(q-1+t)\bigr)\,dt.
\]
由于 \(\tau\) 严格凸，\(\tau'\) 严格单调递增，故上式严格为正，从而
\[
\Delta b(q+1)>\Delta b(q),
\]
即 \(q\) 是 kink。端点在单边定义下自动计入。证毕。
\end{proof}

\begin{theorem}[连续 Legendre 支撑点的一格影随]\label{thm:conclusion-continuous-legendre-support-one-cell-shadowing}
\leanverified{paper\_conclusion\_continuous\_legendre\_support\_one\_cell\_shadowing}
设 \(\tau\) 在区间 \(J\) 上为 \(C^1\) 且严格凸，并令
\[
q_{\mathrm{cont}}(\gamma):=(\tau')^{-1}(\gamma)
\]
为连续 Legendre 支撑点。另一方面，在整数窗 \(\{2,\dots,Q_{\max}\}\) 上定义离散对偶目标
\[
F_\gamma(q):=(q-1)\gamma-\tau(q).
\]
若 \(q_{\mathrm{disc}}(\gamma)\) 是其某个内部最优点，即
\[
q_{\mathrm{disc}}(\gamma)\in \operatorname*{argmax}_{2\le q\le Q_{\max}}F_\gamma(q),
\qquad
3\le q_{\mathrm{disc}}(\gamma)\le Q_{\max}-1,
\]
则必有
\[
\Delta\tau\bigl(q_{\mathrm{disc}}(\gamma)\bigr)
\le \gamma \le
\Delta\tau\bigl(q_{\mathrm{disc}}(\gamma)+1\bigr),
\]
从而
\[
\bigl|q_{\mathrm{cont}}(\gamma)-q_{\mathrm{disc}}(\gamma)\bigr|\le 1.
\]
因此，离散 \(q\)-相图是连续 Legendre 前沿的一格精度整数影随。
\end{theorem}
\begin{proof}
由定理 \ref{thm:conclusion-discrete-convex-budget-kink-fan}，
\[
\Delta\tau\bigl(q_{\mathrm{disc}}(\gamma)\bigr)
\le \gamma \le
\Delta\tau\bigl(q_{\mathrm{disc}}(\gamma)+1\bigr).
\]
由中值定理，存在
\[
\xi_-\in(q_{\mathrm{disc}}(\gamma)-1,q_{\mathrm{disc}}(\gamma)),
\qquad
\xi_+\in(q_{\mathrm{disc}}(\gamma),q_{\mathrm{disc}}(\gamma)+1)
\]
使
\[
\Delta\tau\bigl(q_{\mathrm{disc}}(\gamma)\bigr)=\tau'(\xi_-),
\qquad
\Delta\tau\bigl(q_{\mathrm{disc}}(\gamma)+1\bigr)=\tau'(\xi_+).
\]
故
\[
\tau'(\xi_-)\le \gamma\le \tau'(\xi_+).
\]
由于 \(\tau'\) 严格单调，逆函数存在，从而
\[
q_{\mathrm{cont}}(\gamma)\in[\xi_-,\xi_+]
\subset
[q_{\mathrm{disc}}(\gamma)-1,q_{\mathrm{disc}}(\gamma)+1].
\]
证毕。
\end{proof}

\begin{theorem}[\(q\)-共存平台与离散速率前沿角点]\label{thm:conclusion-q-coexistence-platform-discrete-rate-corner}
设
\[
\Phi(q):=\tau(q)-\log 2,
\qquad
I_{\mathrm{disc}}(\gamma):=\sup_{2\le q\le Q_{\max}}\bigl((q-1)\gamma-\Phi(q)\bigr).
\]
对固定 \(\gamma\)，记
\[
\mathcal M_\gamma:=\operatorname*{argmax}_{2\le q\le Q_{\max}}\bigl((q-1)\gamma-\Phi(q)\bigr).
\]
则 \(\mathcal M_\gamma\) 必为一个连续整数区间。若其非退化，即存在
\[
q_-<q_+,\qquad q_-,q_+\in \mathcal M_\gamma,
\]
则对所有 \(n\in\{q_-,q_-+1,\dots,q_+\}\) 都有
\[
\Phi(n)=\Phi(q_-)+(n-q_-)\gamma.
\]
反之，这种 \(q\)-共存平台当且仅当离散 Legendre 前沿 \(I_{\mathrm{disc}}\) 在 \(\gamma\) 处出现角点；其子梯度恰包含整个区间
\[
[q_--1,q_+-1].
\]
\end{theorem}
\begin{proof}
由定理 \ref{thm:conclusion-discrete-convex-budget-kink-fan}，\(\mathcal M_\gamma\) 由满足相邻斜率夹持条件的整数构成，因此必为连续整数区间。

若 \(q_-<q_+\) 同时最优，则区间中的每个整数都最优。对每个 \(n=q_-,\dots,q_+-1\)，由相邻两点同值，
\[
0=F_\gamma(n+1)-F_\gamma(n)=\gamma-\bigl(\Phi(n+1)-\Phi(n)\bigr),
\]
故
\[
\Phi(n+1)-\Phi(n)=\gamma.
\]
逐项累加即得
\[
\Phi(n)=\Phi(q_-)+(n-q_-)\gamma
\qquad
(q_-\le n\le q_+).
\]

另一方面，\(I_{\mathrm{disc}}\) 是仿射函数族
\[
\gamma\longmapsto (q-1)\gamma-\Phi(q)
\]
的上包络。若至少两条不同斜率的仿射函数在 \(\gamma\) 处同时达到上确界，则 \(I_{\mathrm{disc}}\) 在该点不可微，且其次梯度等于所有活动斜率的凸包，故包含 \([q_--1,q_+-1]\)。反之，若 \(I_{\mathrm{disc}}\) 在 \(\gamma\) 处不可微，则至少两条不同斜率的支撑线同时活动，从而 \(\mathcal M_\gamma\) 非单点。证毕。
\end{proof}

\begin{theorem}[residue 压降的 escort--Rényi 自由能恒等式]\label{thm:conclusion-residue-pressure-drop-escort-renyi-free-energy}
沿定义 \ref{def:pom-residue-refined-moment} 的记号。对每个 \(x\in X_m\) 定义局部桶分布
\[
p_x(i):=\frac{d_{m,i}(x)}{d_m(x)},
\qquad
i=1,\dots,P,
\]
并令
\[
\nu_{m,q}(x):=\frac{d_m(x)^q}{S_q(m)}
\]
为 \(q\)-escort 测度。对 \(q>1\) 的局部 Rényi 熵
\[
H_q(p):=\frac{1}{1-q}\log\sum_{i=1}^{P}p_i^q
\]
有精确恒等式
\[
\frac{S_q^{(P)}(m)}{S_q(m)}
=
\mathbb E_{\nu_{m,q}}\!\left[\sum_{i=1}^{P}p_X(i)^q\right]
=
\mathbb E_{\nu_{m,q}}\!\left[e^{-(q-1)H_q(p_X)}\right].
\]
因此有限尺度压降
\[
\delta_{m,q}^{(P)}:=
\frac1m\log\frac{S_q(m)}{S_q^{(P)}(m)}
\]
满足
\[
\delta_{m,q}^{(P)}
=
-\frac1m\log
\mathbb E_{\nu_{m,q}}\!\left[e^{-(q-1)H_q(p_X)}\right].
\]
也就是说，residue 细分的全部 \(q\)-压降精确等于一项由 escort 取样的局部 Rényi 熵自由能。
\end{theorem}
\begin{proof}
直接计算，
\[
S_q^{(P)}(m)
=
\sum_x\sum_{i=1}^{P}d_{m,i}(x)^q
=
\sum_x d_m(x)^q\sum_{i=1}^{P}\left(\frac{d_{m,i}(x)}{d_m(x)}\right)^q
=
S_q(m)\,
\mathbb E_{\nu_{m,q}}\!\left[\sum_{i=1}^{P}p_X(i)^q\right].
\]
另一方面，
\[
\sum_{i=1}^{P}p_i^q=e^{-(q-1)H_q(p)},
\]
故第二个恒等式成立。再取对数并除以 \(m\)，便得到 \(\delta_{m,q}^{(P)}\) 的自由能表示。证毕。
\end{proof}

\begin{theorem}[二阶细分碰撞的加权均匀化恒等式]\label{thm:conclusion-second-order-refined-collision-weighted-uniformization}
在结论 \ref{thm:conclusion-residue-pressure-drop-escort-renyi-free-energy} 的设定下，进一步假设细分标签来自有限阿贝尔群 \(G\)，并记
\[
p_x(g):=\frac{d_m(x,g)}{d_m(x)},
\qquad
u_G(g):=\frac1{|G|},
\qquad
\nu_{m,2}(x):=\frac{d_m(x)^2}{S_2(m)}.
\]
则二阶细分碰撞满足精确恒等式
\[
\frac{S_2^{(G)}(m)}{S_2(m)}-\frac1{|G|}
=
\mathbb E_{\nu_{m,2}}\|p_X-u_G\|_2^2.
\]
若再记局部非平凡角色能量
\[
\mathcal E_x:=\sum_{\chi\neq\chi_0}|\mu_\chi(x)|^2,
\qquad
\mu_\chi(x):=\sum_{g\in G}p_x(g)\chi(g),
\]
则同一量又可写为
\[
\frac{S_2^{(G)}(m)}{S_2(m)}-\frac1{|G|}
=
\frac1{|G|}\,\mathbb E_{\nu_{m,2}}[\mathcal E_X].
\]
因此，细分碰撞高出 Jensen 地板 \(|G|^{-1}\) 的全部量，恰等于 collision-escort 加权下的局部非均匀性平方，也等于同一加权下的非平凡 Fourier 能量。
\end{theorem}
\begin{proof}
由结论 \ref{thm:conclusion-residue-pressure-drop-escort-renyi-free-energy} 在 \(q=2\) 时，
\[
\frac{S_2^{(G)}(m)}{S_2(m)}
=
\mathbb E_{\nu_{m,2}}\!\left[\sum_{g\in G}p_X(g)^2\right].
\]
而
\[
\|p-u_G\|_2^2
=
\sum_{g\in G}\left(p(g)-\frac1{|G|}\right)^2
=
\sum_{g\in G}p(g)^2-\frac1{|G|},
\]
故第一式成立。

另一方面，由定理 \ref{thm:conclusion-second-refined-collision-exact-fourier-energy}，
\[
\mathcal E_x
=
|G|\sum_{g\in G}p_x(g)^2-1
=
|G|\,\|p_x-u_G\|_2^2.
\]
对 \(\nu_{m,2}\) 取期望即得第二式。证毕。
\end{proof}

\begin{theorem}[局部熵税压力表象与单调比较]\label{thm:conclusion-local-entropy-tax-pressure-representation}
在结论 \ref{thm:conclusion-residue-pressure-drop-escort-renyi-free-energy} 的设定下，对任意 \(q>1\) 定义局部熵税势
\[
\Psi_{m,q}(x):=
q\log d_m(x)-(q-1)H_q(p_x).
\]
则 refined \(q\)-矩具有精确热力学表示
\[
S_q^{(P)}(m)=\sum_{x\in X_m}e^{\Psi_{m,q}(x)}.
\]
因此 refined 有限尺度压力
\[
P_m^{(P)}(q):=\frac1m\log S_q^{(P)}(m)-q\log 2
\]
可写为
\[
P_m^{(P)}(q)
=
\frac1m\log\sum_{x\in X_m}e^{\,q\log d_m(x)-(q-1)H_q(p_x)}
-q\log 2.
\]
特别地，若两种细分方案 \(A,B\) 具有相同的基底纤维计数 \(d_m(x)\)，且对所有 \(x\) 都满足
\[
H_q(p_x^{A})\le H_q(p_x^{B}),
\]
则
\[
P_{m,A}^{(P)}(q)\ge P_{m,B}^{(P)}(q).
\]
因此，在固定 base fold 上，较小的局部 Rényi 熵税严格支配较大的 refined \(q\)-压力。
\end{theorem}
\begin{proof}
由结论 \ref{thm:conclusion-residue-pressure-drop-escort-renyi-free-energy}，
\[
S_q^{(P)}(m)
=
\sum_x d_m(x)^q e^{-(q-1)H_q(p_x)}
=
\sum_x e^{\,q\log d_m(x)-(q-1)H_q(p_x)}.
\]
于是压力表示立即成立。

若 \(A,B\) 具有相同的 \(d_m(x)\) 且
\[
H_q(p_x^{A})\le H_q(p_x^{B})
\qquad (\forall x),
\]
则逐点有
\[
q\log d_m(x)-(q-1)H_q(p_x^{A})
\ge
q\log d_m(x)-(q-1)H_q(p_x^{B}).
\]
对两边取指数、求和并取对数，即得
\[
P_{m,A}^{(P)}(q)\ge P_{m,B}^{(P)}(q).
\]
证毕。
\end{proof}

\begin{theorem}[连续--离散 Legendre 前沿的 Bregman 格点量化]\label{thm:conclusion-continuous-discrete-legendre-bregman-quantization}
\leanverified{paper\_conclusion\_continuous\_discrete\_legendre\_bregman\_quantization}
设 \(\Theta\subset[0,\infty)\) 为闭区间，\(\Lambda\in C^1(\Theta^\circ)\) 且严格凸。对任意
\[
\gamma\in \Lambda'(\Theta^\circ),
\]
记 \(t_\ast(\gamma)\in\Theta^\circ\) 为唯一切触点，即
\[
\Lambda'(t_\ast(\gamma))=\gamma.
\]
定义连续 Legendre 前沿与其整数格近似为
\[
I_\Theta(\gamma):=\sup_{t\in\Theta}\bigl(t\gamma-\Lambda(t)\bigr),
\qquad
I_{\Theta,\ZZ}(\gamma):=\sup_{n\in\Theta\cap\ZZ}\bigl(n\gamma-\Lambda(n)\bigr).
\]
则有精确恒等式
\[
I_\Theta(\gamma)-I_{\Theta,\ZZ}(\gamma)
=
\min_{n\in\Theta\cap\ZZ}
D_\Lambda\bigl(n,t_\ast(\gamma)\bigr),
\]
其中
\[
D_\Lambda(x,y):=\Lambda(x)-\Lambda(y)-\Lambda'(y)(x-y)
\]
为 \(\Lambda\) 的 Bregman 散度。更强地，任一整数最优点
\[
n_\ast(\gamma)\in\operatorname*{argmax}_{n\in\Theta\cap\ZZ}\bigl(n\gamma-\Lambda(n)\bigr)
\]
都满足
\[
n_\ast(\gamma)\in
\{\lfloor t_\ast(\gamma)\rfloor,\ \lceil t_\ast(\gamma)\rceil\}\cap\Theta.
\]
若进一步 \(\Lambda\in C^2(\Theta^\circ)\)，且
\[
\operatorname{dist}\bigl(t_\ast(\gamma),\partial\Theta\bigr)\ge \frac12,
\]
则
\[
0\le I_\Theta(\gamma)-I_{\Theta,\ZZ}(\gamma)
\le
\frac18\sup_{|t-t_\ast(\gamma)|\le 1/2}\Lambda''(t).
\]
\end{theorem}
\begin{proof}
由严格凸性，函数
\[
G_\gamma(t):=t\gamma-\Lambda(t)
\]
严格凹，且其唯一极大点为 \(t=t_\ast(\gamma)\)。因此
\[
I_\Theta(\gamma)=t_\ast(\gamma)\gamma-\Lambda(t_\ast(\gamma)).
\]
对任意整数 \(n\in\Theta\cap\ZZ\)，有
\[
\begin{aligned}
I_\Theta(\gamma)-\bigl(n\gamma-\Lambda(n)\bigr)
&=
\Lambda(n)-\Lambda(t_\ast)-\gamma(n-t_\ast)\\
&=
\Lambda(n)-\Lambda(t_\ast)-\Lambda'(t_\ast)(n-t_\ast)\\
&=
D_\Lambda(n,t_\ast),
\end{aligned}
\]
从而对整数 \(n\) 取下确界便得到
\[
I_\Theta(\gamma)-I_{\Theta,\ZZ}(\gamma)
=
\min_{n\in\Theta\cap\ZZ}D_\Lambda(n,t_\ast(\gamma)).
\]

再证最近格点局域化。由于 \(G_\gamma\) 在 \(t_\ast\) 左侧严格递增、右侧严格递减，故所有严格小于 \(\lfloor t_\ast\rfloor\) 的整数值都不可能优于 \(\lfloor t_\ast\rfloor\)，所有严格大于 \(\lceil t_\ast\rceil\) 的整数值都不可能优于 \(\lceil t_\ast\rceil\)。于是任一整数最优点只能落在这两个最近格点之一。

若 \(\Lambda\in C^2\) 且 \(\operatorname{dist}(t_\ast,\partial\Theta)\ge 1/2\)，取最近整数 \(n^\sharp\) 使
\[
|n^\sharp-t_\ast|\le \frac12.
\]
由二阶 Taylor 公式，存在 \(\xi\) 介于 \(n^\sharp\) 与 \(t_\ast\) 之间，使
\[
D_\Lambda(n^\sharp,t_\ast)=\frac12\Lambda''(\xi)(n^\sharp-t_\ast)^2.
\]
故
\[
I_\Theta(\gamma)-I_{\Theta,\ZZ}(\gamma)
\le
D_\Lambda(n^\sharp,t_\ast)
\le
\frac12\sup_{|t-t_\ast|\le 1/2}\Lambda''(t)\cdot \frac14,
\]
即得所示上界。证毕。
\end{proof}

\begin{theorem}[离散相带宽度的曲率采样]\label{thm:conclusion-discrete-phase-band-width-curvature-sampling}
\leanverified{paper\_conclusion\_discrete\_phase\_band\_width\_curvature\_sampling}
设 \(\Lambda\) 在包含整数窗 \(\{1,\dots,Q_{\max}-1\}\) 的开区间上属于 \(C^2\)，并对每个内点 \(q\in\{3,\dots,Q_{\max}-1\}\) 定义离散相带
\[
\Gamma_q:=
\bigl[
\Lambda(q-1)-\Lambda(q-2),\
\Lambda(q)-\Lambda(q-1)
\bigr].
\]
则：
\begin{enumerate}
  \item \(q\) 成为离散 Legendre 前沿的最优阶，当且仅当
  \[
  \gamma\in \Gamma_q;
  \]
  \item 相带宽度满足精确恒等式
  \[
  |\Gamma_q|
  =
  \Lambda(q)-2\Lambda(q-1)+\Lambda(q-2);
  \]
  \item 存在某个 \(\xi_q\in(q-2,q)\) 使
  \[
  |\Gamma_q|=\Lambda''(\xi_q).
  \]
\end{enumerate}
因此，每一条 \(q\)-相带的宽度恰是连续压力曲线在长度 \(2\) 的窗口内留下的一枚曲率样本。
\end{theorem}
\begin{proof}
第一条是定理 \ref{thm:conclusion-discrete-convex-budget-kink-fan} 应用于离散预算
\[
b(q):=\Lambda(q-1)
\]
后的直接结论。第二条由端点相减立得：
\[
|\Gamma_q|
=
\bigl[\Lambda(q)-\Lambda(q-1)\bigr]
-
\bigl[\Lambda(q-1)-\Lambda(q-2)\bigr]
=
\Lambda(q)-2\Lambda(q-1)+\Lambda(q-2).
\]
第三条则由二阶有限差分的拉格朗日型中值公式给出：对 \(f=\Lambda\) 与步长 \(1\)，存在 \(\xi_q\in(q-2,q)\) 使
\[
\Lambda(q)-2\Lambda(q-1)+\Lambda(q-2)=\Lambda''(\xi_q).
\]
证毕。
\end{proof}

\begin{theorem}[离散扇宽序列的帐篷核层析]\label{thm:conclusion-discrete-fan-width-tent-kernel-tomography}
\leanverified{paper\_conclusion\_discrete\_fan\_width\_tent\_kernel\_tomography}
设 \(\Lambda:\Theta\to\RR\) 仅要求凸，并记其分布意义下的二阶导数为正 Radon 测度
\[
\mu_\Lambda:=\Lambda''.
\]
对每个整数 \(q\) 定义相带宽度
\[
w_q:=\Lambda(q)-2\Lambda(q-1)+\Lambda(q-2).
\]
则有精确表示
\[
w_q
=
\int_{q-2}^{q}
\Bigl(1-\bigl|t-(q-1)\bigr|\Bigr)\,d\mu_\Lambda(t).
\]
若进一步 \(\Lambda\in C^2\)，则
\[
w_q
=
\int_{q-2}^{q}
\Bigl(1-\bigl|t-(q-1)\bigr|\Bigr)\Lambda''(t)\,dt.
\]
因此，离散相带宽度序列 \((w_q)\) 恰是连续曲率测度 \(\mu_\Lambda\) 经帐篷核
\[
K(u):=(1-|u|)_+
\]
卷积后在整数格上的采样。
\end{theorem}
\begin{proof}
先设 \(\Lambda\in C^2\)。对任意 \(x\in\Theta^\circ\)，标准积分恒等式给出
\[
\Lambda(x+1)-2\Lambda(x)+\Lambda(x-1)
=
\int_{-1}^{1}(1-|u|)\Lambda''(x+u)\,du.
\]
令 \(x=q-1\)，并作变量替换 \(t=x+u\)，得到
\[
w_q
=
\int_{q-2}^{q}
\Bigl(1-\bigl|t-(q-1)\bigr|\Bigr)\Lambda''(t)\,dt.
\]

若 \(\Lambda\) 仅凸，则其一阶导数为有界变差函数，分布导数 \(\mu_\Lambda=d\Lambda'\) 为正测度。取标准 mollifier \(\rho_\varepsilon\)，令
\[
\Lambda_\varepsilon:=\rho_\varepsilon\ast\Lambda.
\]
则 \(\Lambda_\varepsilon\in C^\infty\)、\(\Lambda_\varepsilon\to\Lambda\) 局部一致，且 \(\Lambda_\varepsilon''\,dt\rightharpoonup \mu_\Lambda\) 弱收敛。对 \(\Lambda_\varepsilon\) 应用上述光滑公式并令 \(\varepsilon\to 0\)，便得到测度版本。证毕。
\end{proof}

\begin{theorem}[输出分桶几何的扇保持性]\label{thm:conclusion-output-bucket-geometry-preserves-discrete-fan}
\leanverified{paper\_conclusion\_output\_bucket\_geometry\_preserves\_discrete\_fan}
设 \(T,T'\) 为两套确定性在线核，具有同一状态集与同一输出字母表。若对每个输出符号 \(b\in A_{\mathrm{out}}\) 都有
\[
B_b(T)=B_b(T'),
\]
则：
\begin{enumerate}
  \item 对所有实参数 \(t\ge 0\)，两者的射影转移算子完全一致，
  \[
  \mathcal L_{t+1}^{(T)}=\mathcal L_{t+1}^{(T')};
  \]
  \item 因而其连续压力曲线完全一致，
  \[
  \Lambda_T(t)=\Lambda_{T'}(t)
  \qquad(\forall t\ge 0);
  \]
  \item 特别地，整数矩谱塔、离散 Legendre 前沿、相带 \(\Gamma_q\) 与扇宽 \(w_q\) 全部一致。
\end{enumerate}
因此，离散 \(q\)-相图与其连续压力几何首先是输出分桶几何的不变量，而不是输入实现细节的不变量。
\end{theorem}
\begin{proof}
由定义 \ref{def:pom-projective-transfer-operator}，
\[
g_b(w)=\langle\mathbf 1,B_bw\rangle,
\qquad
\Phi_b(w)=\frac{B_bw}{\langle\mathbf 1,B_bw\rangle}
\]
完全由 \(B_b\) 决定。因此若 \(T,T'\) 的全部 \(B_b\) 相同，则对任意实参数 \(q=t+1\ge 1\) 都有
\[
\mathcal L_q^{(T)}f
=
\sum_{b\in A_{\mathrm{out}}}g_b(w)^q f(\Phi_b(w))
=
\mathcal L_q^{(T')}f
\]
作为作用在同一函数空间上的算子恒等相同。于是 Perron 主特征值相同，故
\[
\lambda_T(q)=\lambda_{T'}(q)
\qquad (\forall q\ge 1).
\]
按注 \ref{rem:pom-projective-pressure-route} 的定义
\[
\Lambda(t)=\log\lambda(t+1)-\log|A_{\mathrm{out}}|,
\]
便得到
\[
\Lambda_T(t)=\Lambda_{T'}(t)
\qquad (\forall t\ge 0).
\]
整数矩谱塔 \((r_q)\) 与 \(\Lambda(q-1)\) 的一致性由定理 \ref{thm:pom-projective-operator-degenerates-to-moment-kernel} 给出，而 \(I_{\mathrm{disc}},\Gamma_q,w_q\) 仅由整数锚点 \(\Lambda(q-1)\) 与其有限差分确定，故全部一致。证毕。
\end{proof}

\begin{theorem}[仿射闭式的扇宽排斥]\label{thm:conclusion-affine-closed-form-excluded-by-fan-width}
\leanverified{paper\_conclusion\_affine\_closed\_form\_excluded\_by\_fan\_width}
设在某个区间 \(\Theta\subset[1,\infty)\) 上存在闭式
\[
\lambda(t)=2^t\varphi^{\,1-t}.
\]
令对应连续压力定义为
\[
\Lambda(s):=\log\lambda(s+1)-\log 2
\qquad (s\in\Theta-1).
\]
则
\[
\Lambda(s)=s\log\frac{2}{\varphi},
\]
从而对所有满足 \(q-2,q-1,q\in\Theta-1\) 的整数 \(q\ge 3\)，都有
\[
\Gamma_q=
\left\{\log\frac{2}{\varphi}\right\},
\qquad
w_q=0.
\]
因此，只要在某个整数窗上观测到
\[
w_q>0,
\]
则上述仿射闭式不可能在该窗上成立。
\end{theorem}
\begin{proof}
由假设，
\[
\lambda(s+1)=2^{s+1}\varphi^{-s}.
\]
代入连续压力定义，得到
\[
\Lambda(s)
=
\log\lambda(s+1)-\log 2
=
\log(2^{s+1}\varphi^{-s})-\log 2
=
s\log\frac{2}{\varphi}.
\]
故 \(\Lambda\) 仿射，其一阶差分恒为常数
\[
\Lambda(q)-\Lambda(q-1)=\log\frac{2}{\varphi}.
\]
于是
\[
\Gamma_q=
\Bigl[
\log\frac{2}{\varphi},
\log\frac{2}{\varphi}
\Bigr]
=
\left\{\log\frac{2}{\varphi}\right\},
\]
并且
\[
w_q=\Lambda(q)-2\Lambda(q-1)+\Lambda(q-2)=0.
\]
证毕。
\end{proof}

\begin{theorem}[单元曲率的二矩完备性]\label{thm:conclusion-cell-curvature-two-moment-completeness}
\leanverified{paper\_conclusion\_cell\_curvature\_two\_moment\_completeness}
设 \(n\ge 1\) 为整数，并假设
\[
\operatorname{supp}\mu_\Lambda\subset [n,n+1].
\]
则恰有两条离散扇宽可能非零，即 \(w_{n+1}\) 与 \(w_{n+2}\)，并且它们精确给出该单元内曲率测度的总质量与一阶矩：
\[
w_{n+1}=\int_{[n,n+1]}(n+1-t)\,d\mu_\Lambda(t),
\qquad
w_{n+2}=\int_{[n,n+1]}(t-n)\,d\mu_\Lambda(t).
\]
从而
\[
w_{n+1}+w_{n+2}=\mu_\Lambda([n,n+1]),
\]
以及
\[
n\,w_{n+1}+(n+1)\,w_{n+2}
=
\int_{[n,n+1]} t\,d\mu_\Lambda(t).
\]
\end{theorem}
\begin{proof}
由定理 \ref{thm:conclusion-discrete-fan-width-tent-kernel-tomography}，
\[
w_q=\int_\Theta K\bigl(t-(q-1)\bigr)\,d\mu_\Lambda(t),
\qquad
K(u)=(1-|u|)_+.
\]
由于 \(K(\cdot-(q-1))\) 的支撑为 \([q-2,q]\)，若 \(\operatorname{supp}\mu_\Lambda\subset[n,n+1]\)，则只有 \(q=n+1,n+2\) 的帐篷核与该单元相交。并且在 \([n,n+1]\) 上，
\[
K(t-n)=n+1-t,
\qquad
K\bigl(t-(n+1)\bigr)=t-n.
\]
代入即可得到前两式。

两式相加给出
\[
w_{n+1}+w_{n+2}
=
\int_{[n,n+1]}\bigl((n+1-t)+(t-n)\bigr)\,d\mu_\Lambda(t)
=
\mu_\Lambda([n,n+1]).
\]
再作线性组合
\[
n(n+1-t)+(n+1)(t-n)=t,
\]
对 \(\mu_\Lambda\) 积分即得一阶矩公式。证毕。
\end{proof}

\begin{corollary}[单元内曲率重心的邻峰比值公式]\label{cor:conclusion-cell-curvature-barycenter-adjacent-peak-ratio}
\leanverified{paper\_conclusion\_cell\_curvature\_barycenter\_adjacent\_peak\_ratio}
在定理 \ref{thm:conclusion-cell-curvature-two-moment-completeness} 的设定下，若
\[
\mu_\Lambda([n,n+1])>0,
\]
则该单元内曲率测度的重心
\[
\bar t_n:=
\frac{\int_{[n,n+1]} t\,d\mu_\Lambda(t)}
{\mu_\Lambda([n,n+1])}
\]
可由离散扇宽精确恢复为
\[
\bar t_n
=
n+\frac{w_{n+2}}{w_{n+1}+w_{n+2}}.
\]
\end{corollary}
\begin{proof}
由定理 \ref{thm:conclusion-cell-curvature-two-moment-completeness}，
\[
\int_{[n,n+1]} t\,d\mu_\Lambda(t)
=
n\,w_{n+1}+(n+1)\,w_{n+2},
\]
且
\[
\mu_\Lambda([n,n+1])=w_{n+1}+w_{n+2}.
\]
两式相除即可。证毕。
\end{proof}

\begin{theorem}[单隐藏相变点的精确反演]\label{thm:conclusion-single-hidden-phase-transition-exact-inversion}
\leanverified{paper\_conclusion\_single\_hidden\_phase\_transition\_exact\_inversion}
沿用定理 \ref{thm:conclusion-cell-curvature-two-moment-completeness} 的设定。若更强地
\[
\mu_\Lambda=a\,\delta_\tau,
\qquad
\tau\in[n,n+1],\quad a>0,
\]
则离散扇宽满足
\[
w_{n+1}=a(n+1-\tau),
\qquad
w_{n+2}=a(\tau-n).
\]
因而该隐藏相变点的强度与位置可显式反演为
\[
a=w_{n+1}+w_{n+2},
\qquad
\tau=n+\frac{w_{n+2}}{w_{n+1}+w_{n+2}}.
\]
\end{theorem}
\begin{proof}
把 \(\mu_\Lambda=a\delta_\tau\) 代入定理 \ref{thm:conclusion-cell-curvature-two-moment-completeness} 的两条公式，立即得到
\[
w_{n+1}=a(n+1-\tau),
\qquad
w_{n+2}=a(\tau-n).
\]
两式相加给出 \(a\)；再由第二式除以和式得到 \(\tau\)。证毕。
\end{proof}

\begin{theorem}[尾曲率预算与连续冻结的等价]\label{thm:conclusion-tail-curvature-budget-equivalent-continuous-freezing}
对任意整数 \(N\ge 2\)，定义尾部权函数
\[
\eta_N(t):=\sum_{q=N}^{\infty}K\bigl(t-(q-1)\bigr).
\]
则有精确闭式
\[
\eta_N(t)=
\begin{cases}
0, & t\le N-2,\\[2mm]
t-(N-2), & N-2\le t\le N-1,\\[2mm]
1, & t\ge N-1.
\end{cases}
\]
从而
\[
\sum_{q=N}^{\infty}w_q
=
\int_{N-2}^{N-1}\bigl(t-(N-2)\bigr)\,d\mu_\Lambda(t)
\;+\;
\mu_\Lambda([N-1,\infty)).
\]
特别地，成立严格等价：
\[
w_q=0\quad(\forall q\ge N)
\iff
\mu_\Lambda\bigl((N-2,\infty)\bigr)=0.
\]
等价地，\(\Lambda\) 在开半轴 \((N-2,\infty)\) 上恰为仿射函数。若再有 \(\mu_\Lambda(\{N-2\})=0\)，则 \(\Lambda\) 在闭半轴 \([N-2,\infty)\) 上仿射。
\leanverified{paper\_conclusion\_tail\_curvature\_budget\_equivalent\_continuous\_freezing}
\end{theorem}
\begin{proof}
对 \(t\le N-2\)，所有中心 \(q-1\ge N-1\) 与 \(t\) 的距离都至少为 \(1\)，故 \(\eta_N(t)=0\)。

若 \(t\in[N-2,N-1]\)，则只有 \(q=N\) 这一项非零，于是
\[
\eta_N(t)=K\bigl(t-(N-1)\bigr)=t-(N-2).
\]

若 \(t\ge N-1\)，则所有可能非零的帐篷核中心都已落在求和范围之内；由平移帐篷核的分片线性分割恒等式，
\[
\sum_{m\in\ZZ}K(t-m)=1,
\]
且对给定 \(t\ge N-1\) 的非零项都满足 \(m\ge N-1\)，故 \(\eta_N(t)=1\)。

再由 Tonelli 定理交换求和与积分，
\[
\sum_{q=N}^{\infty}w_q
=
\int \eta_N(t)\,d\mu_\Lambda(t),
\]
代入上述闭式便得尾预算公式。

由于 \(\eta_N(t)\ge 0\)，且在 \((N-2,\infty)\) 上严格为正，故
\[
\sum_{q=N}^{\infty}w_q=0
\iff
\mu_\Lambda\bigl((N-2,\infty)\bigr)=0.
\]
又因 \(w_q\ge 0\)，和为零当且仅当每一项都为零，即
\[
\sum_{q=N}^{\infty}w_q=0
\iff
w_q=0\quad(\forall q\ge N).
\]
最后，凸函数在某开区间上二阶分布导数为零当且仅当其在该区间上仿射；若端点 \(N-2\) 也无原子，则仿射性可延伸到闭半轴。证毕。
\end{proof}

\begin{theorem}[孤立扇峰的整数刚性]\label{thm:conclusion-isolated-fan-peak-integer-rigidity}
\leanverified{paper\_conclusion\_isolated\_fan\_peak\_integer\_rigidity}
设 \(q\ge 3\)。若出现一个孤立非零扇峰
\[
w_{q-1}=0,\qquad
w_q>0,\qquad
w_{q+1}=0,
\]
则在局部窗口 \([q-2,q]\) 内，曲率测度被强制坍缩为单点原子：
\[
\mu_\Lambda\bigl([q-2,q]\setminus\{q-1\}\bigr)=0,
\qquad
\mu_\Lambda(\{q-1\})=w_q.
\]
因此 \(t=q-1\) 是一个精确整数相变点，其跳跃质量正好等于该孤立扇峰的高度。
\end{theorem}
\begin{proof}
对 \(t\in[q-2,q-1)\)，有
\[
K\bigl(t-(q-2)\bigr)=q-1-t>0.
\]
若该区间上有正质量，则
\[
w_{q-1}=\int K\bigl(t-(q-2)\bigr)\,d\mu_\Lambda(t)>0,
\]
与 \(w_{q-1}=0\) 矛盾。因此
\[
\mu_\Lambda([q-2,q-1))=0.
\]
同理，对 \(t\in(q-1,q]\)，有
\[
K(t-q)=t-q+1>0.
\]
若该区间上有正质量，则
\[
w_{q+1}=\int K(t-q)\,d\mu_\Lambda(t)>0,
\]
与 \(w_{q+1}=0\) 矛盾。故
\[
\mu_\Lambda((q-1,q])=0.
\]
于是 \([q-2,q]\) 内唯一可能的支撑点是 \(q-1\)。因此
\[
w_q=\int K\bigl(t-(q-1)\bigr)\,d\mu_\Lambda(t)=K(0)\mu_\Lambda(\{q-1\})=\mu_\Lambda(\{q-1\}),
\]
因为 \(K(0)=1\)。证毕。
\end{proof}

\begin{theorem}[整数支撑曲率的精确样条重建]\label{thm:conclusion-integer-supported-curvature-exact-spline-reconstruction}
\leanverified{paper\_conclusion\_integer\_supported\_curvature\_exact\_spline\_reconstruction}
假设
\[
\mu_\Lambda((0,1))=0,
\qquad
\operatorname{supp}\mu_\Lambda\subset \ZZ_{\ge 1}.
\]
则存在一列非负系数 \((\kappa_n)_{n\ge 1}\) 使
\[
\mu_\Lambda=\sum_{n\ge 1}\kappa_n\delta_n,
\qquad
\kappa_n=w_{n+1}.
\]
并且整个连续压力曲线可由离散扇宽显式重建为
\[
\Lambda(t)
=
\Lambda(0)+\bigl(\Lambda(1)-\Lambda(0)\bigr)t
\;+\;
\sum_{n\ge 1}w_{n+1}(t-n)_+,
\qquad
t\ge 0.
\]
因此，在整数支撑制度下，离散扇宽序列 \((w_q)_{q\ge 2}\) 与两个端点锚 \(\Lambda(0),\Lambda(1)\) 已足以唯一决定整条连续压力曲线与其 Legendre 前沿。
\end{theorem}
\begin{proof}
由支撑假设，可写
\[
\mu_\Lambda=\sum_{n\ge 1}\kappa_n\delta_n.
\]
于是对每个 \(q\ge 2\)，
\[
w_q
=
\sum_{n\ge 1}\kappa_n K\bigl(n-(q-1)\bigr).
\]
而当 \(n,q\) 为整数时，
\[
K\bigl(n-(q-1)\bigr)=
\begin{cases}
1,& n=q-1,\\
0,& n\neq q-1,
\end{cases}
\]
故
\[
w_q=\kappa_{q-1},
\qquad
\kappa_n=w_{n+1}.
\]

再定义
\[
\Psi(t):=
\Lambda(0)+\bigl(\Lambda(1)-\Lambda(0)\bigr)t
\;+\;
\sum_{n\ge 1}w_{n+1}(t-n)_+.
\]
对每个固定 \(t\)，上式只有有限项非零，故 \(\Psi\) 良定义且为凸分段仿射函数。其分布二阶导数满足
\[
\Psi''=\sum_{n\ge 1}w_{n+1}\delta_n=\mu_\Lambda=\Lambda''.
\]
因此 \(\Lambda-\Psi\) 的二阶分布导数为零，故 \(\Lambda-\Psi\) 为仿射函数。又因 \(\mu_\Lambda((0,1))=0\)，\(\Lambda\) 在 \([0,1]\) 上仿射；而 \(\Psi(0)=\Lambda(0)\)、\(\Psi(1)=\Lambda(1)\)，故这条仿射差在两个点上为零，只能恒为零。于是 \(\Lambda\equiv\Psi\)。证毕。
\end{proof}

\begin{corollary}[孤立峰族的全局 polyhedral spline 闭包]\label{cor:conclusion-isolated-peak-family-global-polyhedral-spline-closure}
\leanverified{paper\_conclusion\_isolated\_peak\_family\_global\_polyhedral\_spline\_closure}
若除边界胞 \((0,1)\) 外，所有正扇峰都满足定理 \ref{thm:conclusion-isolated-fan-peak-integer-rigidity} 的孤立条件，则
\[
\operatorname{supp}\mu_\Lambda\cap[1,\infty)\subset \ZZ_{\ge 1},
\]
从而定理 \ref{thm:conclusion-integer-supported-curvature-exact-spline-reconstruction} 的整数支撑假设在尾部自动成立。于是 \(\Lambda\) 在 \([1,\infty)\) 上成为一个由离散扇峰完全决定的 polyhedral spline。
\end{corollary}
\begin{proof}
若某个整数 \(n\ge 1\) 不是曲率支撑点，则对应 \(w_{n+1}=0\)。若 \(w_{n+1}>0\)，按假设其为孤立扇峰，定理 \ref{thm:conclusion-isolated-fan-peak-integer-rigidity} 蕴含
\[
\mu_\Lambda([n-1,n+1]\setminus\{n\})=0.
\]
因此每个正扇峰只对应一个整数原子点，且尾部不存在非整数支撑。再由定理 \ref{thm:conclusion-integer-supported-curvature-exact-spline-reconstruction} 即得结论。证毕。
\end{proof}

\begin{theorem}[近孤立扇峰的定量局部化]\label{thm:conclusion-near-isolated-fan-peak-quantitative-localization}
\leanverified{paper\_conclusion\_near\_isolated\_fan\_peak\_quantitative\_localization}
设 \(q\ge 3\)、\(\varepsilon\in(0,1)\)。定义局部离心质量
\[
M^-_{q,\varepsilon}:=
\mu_\Lambda\bigl([q-2,\ q-1-\varepsilon]\bigr),
\qquad
M^+_{q,\varepsilon}:=
\mu_\Lambda\bigl([q-1+\varepsilon,\ q]\bigr).
\]
则有精确上界
\[
M^-_{q,\varepsilon}\le \frac{w_{q-1}}{\varepsilon},
\qquad
M^+_{q,\varepsilon}\le \frac{w_{q+1}}{\varepsilon},
\]
从而
\[
\mu_\Lambda\Bigl([q-2,q]\setminus(q-1-\varepsilon,\ q-1+\varepsilon)\Bigr)
\le
\frac{w_{q-1}+w_{q+1}}{\varepsilon}.
\]
\end{theorem}
\begin{proof}
对 \(t\in[q-2,q-1-\varepsilon]\)，有
\[
K\bigl(t-(q-2)\bigr)=q-1-t\ge \varepsilon.
\]
因此
\[
w_{q-1}
=
\int K\bigl(t-(q-2)\bigr)\,d\mu_\Lambda(t)
\ge
\int_{[q-2,q-1-\varepsilon]}\varepsilon\,d\mu_\Lambda(t)
=
\varepsilon M^-_{q,\varepsilon},
\]
故
\[
M^-_{q,\varepsilon}\le \frac{w_{q-1}}{\varepsilon}.
\]
同理，对 \(t\in[q-1+\varepsilon,q]\)，有
\[
K(t-q)=t-q+1\ge \varepsilon,
\]
于是
\[
w_{q+1}\ge \varepsilon M^+_{q,\varepsilon},
\qquad
M^+_{q,\varepsilon}\le \frac{w_{q+1}}{\varepsilon}.
\]
两式相加即得总离心质量上界。证毕。
\end{proof}

\begin{corollary}[近孤立扇峰的弱骨架化]\label{cor:conclusion-near-isolated-fan-peak-weak-skeletonization}
\leanverified{paper\_conclusion\_near\_isolated\_fan\_peak\_weak\_skeletonization}
设 \(q_j\to\infty\) 为一列整数，并记
\[
\nu_j:=
\frac{\mu_\Lambda|_{[q_j-2,q_j]}}{\mu_\Lambda([q_j-2,q_j])},
\qquad
\mu_\Lambda([q_j-2,q_j])>0.
\]
若
\[
\frac{w_{q_j-1}+w_{q_j+1}}{\mu_\Lambda([q_j-2,q_j])}\longrightarrow 0,
\]
则平移后的局部曲率分布
\[
E\longmapsto \nu_j(E+q_j-1)
\]
弱收敛到 \(\delta_0\)。换言之，局部曲率分布骨架化到整数中心 \(q_j-1\)。
\end{corollary}
\begin{proof}
由定理 \ref{thm:conclusion-near-isolated-fan-peak-quantitative-localization}，对任意固定 \(\varepsilon\in(0,1)\)，
\[
\nu_j\Bigl([q_j-2,q_j]\setminus(q_j-1-\varepsilon,q_j-1+\varepsilon)\Bigr)
\le
\frac{w_{q_j-1}+w_{q_j+1}}{\varepsilon\,\mu_\Lambda([q_j-2,q_j])}
\longrightarrow 0.
\]
因此任意 \(\varepsilon\)-邻域外的质量都趋于零，故 \(\nu_j\) 弱收敛到 \(\delta_{q_j-1}\)。平移后即得到 \(\delta_0\)。证毕。
\end{proof}


\begin{conjecture}[帐篷核层析的唯一重建]\label{conj:conclusion-tent-kernel-tomography-unique-reconstruction}
设连续压力 \(\Lambda\) 满足附录 \ref{app:pom-projective-pressure-operator} 所要求的谱隙--解析制度，并且其曲率测度 \(\mu_\Lambda\) 落在一个足够窄的 Paley--Wiener 型带宽类中。则定理 \ref{thm:conclusion-discrete-fan-width-tent-kernel-tomography} 的帐篷层析应是唯一可逆的：离散相带宽度序列
\[
(w_q)_{q\ge 3}
\]
唯一决定 \(\mu_\Lambda\)，从而唯一决定 \(\Lambda\)（差一个仿射项）；而该仿射项又由整点锚定
\[
\Lambda(q-1)=\log r_q-\log|A_{\mathrm{out}}|
\]
固定。等价地，整条连续 Legendre 前沿 \(I\) 应可由离散扇 \((\Gamma_q)_{q\ge 3}\) 唯一恢复。

支持性理由在于：定理 \ref{thm:conclusion-discrete-fan-width-tent-kernel-tomography} 已给出
\[
w_q=(K\ast\mu_\Lambda)(q-1),
\qquad
K(u)=(1-|u|)_+,
\]
而其 Fourier 变换满足
\[
\widehat K(\xi)=\left(\frac{\sin(\xi/2)}{\xi/2}\right)^2.
\]
在低频带内，\(\widehat K\) 无零点。若谱隙解析性把 \(\mu_\Lambda\) 压入这一无零频带，则帐篷核卷积应可唯一反演；结合整数锚点即可消去仿射不确定性。
\end{conjecture}

\begin{theorem}[固定分辨率 Pareto 冷端的对数尖点]\label{thm:conclusion-fixed-resolution-pareto-coldend-log-cusp}
\leanverified{paper\_conclusion\_fixed\_resolution\_pareto\_coldend\_log\_cusp}
沿用定理 \ref{thm:pom-finite-pareto-legendre-curvature} 的记号，并设 \(d_m\) 非常值。记顶端三层多重度与对应基数为
\[
M_1:=M_m,\qquad
M_2:=M_m^{(2)},\qquad
M_3:=M_m^{(3)},
\]
\[
N_1:=\#\{x\in X_m:\ d_m(x)=M_1\},
\qquad
N_2:=\#\{x\in X_m:\ d_m(x)=M_2\},
\]
其中约定若 \(d_m\) 只有两层取值，则 \(M_3:=0\)。再记
\[
\kappa_\ast:=\log M_1,
\qquad
H_\ast:=\log N_1,
\qquad
\Delta_m:=\log\frac{M_1}{M_2}>0,
\qquad
a_m:=\frac{N_2}{N_1}.
\]
则当 \(\kappa\uparrow\kappa_\ast\) 时，固定分辨率 Pareto 上边界满足精确尖点展开
\[
\mathcal H_m(\kappa)
=
H_\ast
+
\frac{\kappa_\ast-\kappa}{\Delta_m}
\left(
\log\frac1{\kappa_\ast-\kappa}
+
1+\log(a_m\Delta_m)
\right)
+
o(\kappa_\ast-\kappa).
\]
因此冷端并非二次光滑闭合，而是一个由第二层高度差 \(\Delta_m\) 与第二层相对简并比 \(a_m\) 共同强制的对数尖点。
\end{theorem}
\begin{proof}
由定理 \ref{thm:pom-finite-pareto-legendre-curvature}，
\[
\kappa=F_m'(q),
\qquad
\mathcal H_m(\kappa)=F_m(q)-qF_m'(q),
\qquad
F_m(q):=\log S_q(m),
\]
且当 \(q\to+\infty\) 时有 \(\kappa\uparrow \kappa_\ast=\log M_1\)。

把幂和按顶端三层分拆，得到
\[
S_q(m)
=
N_1M_1^q+N_2M_2^q+O(M_3^q)
=
N_1M_1^q\Bigl(1+a_m\rho^q+O(\sigma^q)\Bigr),
\]
其中
\[
\rho:=\frac{M_2}{M_1}=e^{-\Delta_m}\in(0,1),
\qquad
\sigma:=\frac{M_3}{M_1}\in[0,\rho).
\]
于是
\[
F_m(q)
=
q\log M_1+\log N_1+\log\!\Bigl(1+a_m\rho^q+O(\sigma^q)\Bigr).
\]
对 \(q\) 求导得到
\[
\kappa_\ast-\kappa
=
-\frac{d}{dq}\log\!\Bigl(1+a_m\rho^q+O(\sigma^q)\Bigr)
=
a_m\Delta_m\,\rho^q+o(\rho^q).
\]
因此
\[
\rho^q=\frac{\kappa_\ast-\kappa}{a_m\Delta_m}+o(\kappa_\ast-\kappa),
\]
并且
\[
q
=
\frac1{\Delta_m}
\left(
\log(a_m\Delta_m)-\log(\kappa_\ast-\kappa)
\right)
+o(1).
\]

另一方面，由
\[
\mathcal H_m(\kappa)=F_m(q)-qF_m'(q)
\]
与上式展开可得
\[
\mathcal H_m(\kappa)
=
\log N_1+a_m(1+q\Delta_m)\rho^q+o(\rho^q).
\]
再将前述 \(\rho^q\) 与 \(q\) 的表达式代入并整理，即得
\[
\mathcal H_m(\kappa)
=
H_\ast
+
\frac{\kappa_\ast-\kappa}{\Delta_m}
\left(
\log\frac1{\kappa_\ast-\kappa}
+
1+\log(a_m\Delta_m)
\right)
+
o(\kappa_\ast-\kappa).
\]
证毕。
\end{proof}

\begin{corollary}[冷端有限曲率延拓不可能与二阶极点律]\label{cor:conclusion-pareto-coldend-curvature-pole}
在定理 \ref{thm:conclusion-fixed-resolution-pareto-coldend-log-cusp} 的设定下，当 \(\kappa\uparrow\kappa_\ast\) 时有
\[
\mathcal H_m'(\kappa)
=
-\frac1{\Delta_m}\log\frac1{\kappa_\ast-\kappa}
+O(1),
\]
以及
\[
\mathcal H_m''(\kappa)
=
-\frac{1}{\Delta_m(\kappa_\ast-\kappa)}
+o\!\left(\frac1{\kappa_\ast-\kappa}\right).
\]
特别地，\(\mathcal H_m\) 不可能在 \(\kappa_\ast\) 处具有有限曲率的 \(C^2\)-延拓；冷端几何必为真正的对数尖点。
\leanverified{paper_conclusion_pareto_coldend_curvature_pole_seeds}
\leanverified{paper_conclusion_pareto_coldend_curvature_pole_package}
\end{corollary}
\begin{proof}
由定理 \ref{thm:pom-finite-pareto-legendre-curvature}，
\[
\mathcal H_m'(\kappa)=-q,
\qquad
\mathcal H_m''(\kappa)=-\frac1{F_m''(q)}.
\]
而在定理 \ref{thm:conclusion-fixed-resolution-pareto-coldend-log-cusp} 的证明中已经得到
\[
\kappa_\ast-\kappa=a_m\Delta_m\rho^q+o(\rho^q),
\qquad
q=
\frac1{\Delta_m}\log\frac1{\kappa_\ast-\kappa}
+O(1).
\]
这立刻给出第一式。再由
\[
F_m''(q)=a_m\Delta_m^2\rho^q+o(\rho^q)
=
\Delta_m(\kappa_\ast-\kappa)+o(\kappa_\ast-\kappa),
\]
代入 \(\mathcal H_m''(\kappa)=-1/F_m''(q)\) 即得第二式。证毕。
\end{proof}

\begin{theorem}[冷端尖点对顶端第二层参数的层析]\label{thm:conclusion-pareto-coldend-two-parameter-tomography}
\leanverified{paper\_conclusion\_pareto\_coldend\_two\_parameter\_tomography}
在定理 \ref{thm:conclusion-fixed-resolution-pareto-coldend-log-cusp} 的设定下，冷端尖点满足
\[
\lim_{\kappa\uparrow\kappa_\ast}
\frac{\mathcal H_m(\kappa)-H_\ast}
{(\kappa_\ast-\kappa)\log\frac1{\kappa_\ast-\kappa}}
=
\frac1{\Delta_m},
\]
并且
\[
\lim_{\kappa\uparrow\kappa_\ast}
\left[
\frac{\Delta_m\bigl(\mathcal H_m(\kappa)-H_\ast\bigr)}{\kappa_\ast-\kappa}
-
\log\frac1{\kappa_\ast-\kappa}
\right]
=
1+\log(a_m\Delta_m).
\]
因此，冷端 Pareto 前沿的局部几何已经唯一恢复
\[
\Delta_m=\log\frac{M_1}{M_2},
\qquad
a_m=\frac{N_2}{N_1}.
\]
\end{theorem}
\begin{proof}
把定理 \ref{thm:conclusion-fixed-resolution-pareto-coldend-log-cusp} 的展开写成
\[
\mathcal H_m(\kappa)-H_\ast
=
\frac{y}{\Delta_m}
\left(
\log\frac1y+1+\log(a_m\Delta_m)
\right)
+o(y),
\qquad
y:=\kappa_\ast-\kappa.
\]
分别除以 \(y\log(1/y)\) 与 \(y\)，便得到两个极限。第一个极限给出 \(\Delta_m^{-1}\)，第二个极限则给出 \(a_m\Delta_m\)。故 \(\Delta_m\) 与 \(a_m\) 可被唯一恢复。证毕。
\end{proof}

\begin{theorem}[右端可见相的支撑函数闭包]\label{thm:conclusion-right-edge-visible-phases-support-function-closure}
\leanverified{paper\_conclusion\_right\_edge\_visible\_phases\_support\_function\_closure}
沿用定理 \ref{thm:pom-spectrum-right-edge-support-function} 的记号，并假设定理 \ref{thm:pom-multiwell-softmax-low-temperature-expansion} 中的右端可见相有限，记为
\[
(\alpha_i,h_i)_{i=1}^{K},
\qquad
\alpha_1=\alpha_\ast>\alpha_2>\cdots>\alpha_K,
\qquad
h_i=f(\alpha_i).
\]
定义冷端缺口函数
\[
R(a):=P_a-a\alpha_\ast
\qquad
(a\in\ZZ_{\ge 1}).
\]
则有精确闭式
\[
R(a)=\max_{1\le i\le K}\bigl\{h_i-a(\alpha_\ast-\alpha_i)\bigr\}.
\]
因此 \(R(a)\) 是一条凸、分段仿射、单调不增的右端支撑函数；其每一段仿射片恰对应一个可见相，而任意两相的交叉尺度由
\[
a_{i\leftrightarrow j}
=
\frac{h_i-h_j}{\alpha_i-\alpha_j}
\]
给出。
\end{theorem}
\begin{proof}
由定理 \ref{thm:pom-spectrum-right-edge-support-function}，
\[
R(a)=\sup_{\alpha\in\Sigma_f}\bigl(f(\alpha)-a(\alpha_\ast-\alpha)\bigr).
\]
而在有限可见相假设下，右端有效支撑恰由有限个点 \((\alpha_i,h_i)\) 给出，故上确界退化为有限条仿射函数
\[
a\longmapsto h_i-a(\alpha_\ast-\alpha_i)
\]
的最大值。其余结论皆为有限条直线取上包络的直接推论。证毕。
\end{proof}

\begin{corollary}[不可见相与整数矩阶的暴露失败]\label{cor:conclusion-nonexposed-phase-fails-integer-domination}
在定理 \ref{thm:conclusion-right-edge-visible-phases-support-function-closure} 的设定下，若某一可见相 \((\alpha_j,h_j)\) 不是点集
\[
\{(\alpha_i,h_i)\}_{i=1}^{K}
\]
上凸包的暴露顶点，则对每个整数 \(a\ge 1\) 都有严格不等式
\[
h_j-a(\alpha_\ast-\alpha_j)<R(a).
\]
因此该相虽然可以在低温 softmax 展开中以指数小项出现，却不可能成为任何整数矩阶的主导相。换言之，整数矩阶真正看见的只是右端谱凸包的整数可见暴露相。
\leanverified{paper\_conclusion\_nonexposed\_phase\_integer\_domination\_seeds}
\leanverified{paper\_conclusion\_nonexposed\_phase\_integer\_domination\_package}
\end{corollary}
\begin{proof}
由定理 \ref{thm:conclusion-right-edge-visible-phases-support-function-closure}，
\[
R(a)=\max_{1\le i\le K}\bigl\{h_i-a(\alpha_\ast-\alpha_i)\bigr\}.
\]
若 \((\alpha_j,h_j)\) 不是上凸包的暴露顶点，则其对应仿射函数
\[
L_j(a):=h_j-a(\alpha_\ast-\alpha_j)
\]
在任何 \(a\ge 1\) 上都不可能成为上包络的支撑线。于是对每个整数 \(a\ge 1\) 都有
\[
L_j(a)<\max_i L_i(a)=R(a).
\]
证毕。
\end{proof}

\begin{theorem}[冷端凸包与大阶差分窗口的 Newton--Prony 等价]\label{thm:conclusion-coldend-convex-hull-difference-window-prony-equivalence}
\leanverified{paper\_conclusion\_coldend\_convex\_hull\_difference\_window\_prony\_equivalence}
沿用定理 \ref{thm:conclusion-right-edge-visible-phases-support-function-closure} 的设定，并记
\[
\mathcal C_{\mathrm{cold}}
:=
\operatorname{conv}^{\uparrow}\{(\alpha_i,h_i)\}_{i=1}^{K}
\]
为右端可见相的上凸包。进一步假设 \(\mathcal C_{\mathrm{cold}}\) 的每个暴露顶点的法锥都与 \(\ZZ_{\ge1}\) 相交。则以下三类数据彼此唯一决定：
\begin{enumerate}
  \item 冷端上凸包 \(\mathcal C_{\mathrm{cold}}\)；
  \item 冷端支撑函数
  \[
  R(a)=P_a-a\alpha_\ast
  \qquad (a\in\ZZ_{\ge1});
  \]
  \item 充分大 \(q\) 处的一段有限差分窗口
  \[
  \bigl\{\Delta P(q+j)\bigr\}_{j=0}^{K_0},
  \qquad
  \Delta P(q):=P(q+1)-P(q),
  \]
  其中 \(K_0\) 仅依赖于可见相个数上界。
\end{enumerate}
换言之，冷端凸包几何与大 \(q\) 差分谱在信息论上完全等价；前者是几何表述，后者是可审计算法表述。
\end{theorem}
\begin{proof}
\textup{(i)\(\Rightarrow\)(ii)} 由定理 \ref{thm:conclusion-right-edge-visible-phases-support-function-closure} 立即得到。

\textup{(ii)\(\Rightarrow\)(i)} 时，\(R(a)\) 是有限条仿射函数的上包络。每个暴露片都可写成
\[
a\longmapsto h_i-a(\alpha_\ast-\alpha_i).
\]
由仿射片的斜率与截距即可恢复相应的 \((\alpha_i,h_i)\)。在附加的整数可见性假设下，每个暴露顶点都在某个整数 \(a\) 上被支撑函数看见，因此 \(R(a)\) 的整数数据唯一恢复整个上凸包 \(\mathcal C_{\mathrm{cold}}\)。

\textup{(iii)\(\Rightarrow\)(i)} 时，定理 \ref{thm:pom-multiwell-softmax-low-temperature-expansion} 给出
\[
P(q)-q\alpha_1-h_1
=
\log\!\left(
1+\sum_{i\ge2}e^{-(\alpha_1-\alpha_i)q+(h_i-h_1)}
\right)
+O(e^{-cq}),
\]
故差分序列 \(\Delta P(q)\) 具有有限指数和的可审计扰动结构。再由定理 \ref{thm:pom-multiwell-top-parameter-identifiability} 的 Prony--Hankel 识别，可由充分大的有限差分窗口恢复全部可见竞争相 \((\alpha_i,h_i)\)，因而恢复其上凸包。

\textup{(i)\(\Rightarrow\)(iii)} 则只需把 \(\mathcal C_{\mathrm{cold}}\) 所对应的暴露相数据代回同一多井低温展开，即得到差分窗口的有限指数模型。证毕。
\end{proof}

\begin{theorem}[冷端支撑函数的有限整阶平台]\label{thm:conclusion-coldend-support-function-finite-integer-plateau}
\leanverified{paper\_conclusion\_coldend\_support\_function\_finite\_integer\_plateau}
在定理 \ref{thm:conclusion-right-edge-visible-phases-support-function-closure} 的设定下，定义
\[
a_{\mathrm{plat}}
:=
\max_{i\ge2}\frac{h_i-h_1}{\alpha_1-\alpha_i},
\qquad
\alpha_1=\alpha_\ast,
\qquad
h_1=h_\ast.
\]
则对每个整数 \(a>a_{\mathrm{plat}}\) 都有精确平台
\[
R(a)=h_1=h_\ast,
\qquad
P_a=a\alpha_\ast+h_\ast.
\]
因此冷端支撑函数在有限整阶之后出现真正的水平平台，而不是仅仅渐近接近。若再施加严格基态间隙并调用定理 \ref{thm:fold-pressure-freezing-threshold}，则这一有限整阶平台正是实参数冻结枝
\[
P_t=t\alpha_\ast+h_\ast
\qquad (t>a_{\mathrm c})
\]
在整数温度轴上的离散影像。
\end{theorem}
\begin{proof}
由定理 \ref{thm:conclusion-right-edge-visible-phases-support-function-closure}，
\[
R(a)=\max\Bigl\{h_1,\ \max_{i\ge2}\bigl[h_i-a(\alpha_1-\alpha_i)\bigr]\Bigr\}.
\]
对每个 \(i\ge2\)，一旦
\[
a>\frac{h_i-h_1}{\alpha_1-\alpha_i},
\]
便有
\[
h_i-a(\alpha_1-\alpha_i)<h_1.
\]
因此当 \(a>a_{\mathrm{plat}}\) 时，所有次级相对应的仿射线都严格落在主相水平线 \(h_1\) 之下，从而
\[
R(a)=h_1.
\]
于是
\[
P_a=a\alpha_1+h_1=a\alpha_\ast+h_\ast.
\]
若进一步满足定理 \ref{thm:fold-pressure-freezing-threshold} 的严格基态间隙假设，则对一切实数 \(t>a_{\mathrm c}\) 都有
\[
P_t=t\alpha_\ast+h_\ast,
\]
故上述整数平台正是该冻结枝在整数温度轴上的离散影像。证毕。
\end{proof}

\begin{theorem}[金属平均 Pareto 前沿的暴露端点结构]\label{thm:conclusion-metallic-pareto-exposed-golden-endpoint}
对金属平均族
\[
\alpha_A=[0;A,A,A,\dots],
\qquad
\lambda_A=\frac{A+\sqrt{A^2+4}}{2},
\qquad
h_Z(A)=\log\frac{A+1}{\lambda_A},
\qquad
\kappa(A)=\frac{A}{\log\lambda_A},
\]
定义线性标量化目标
\[
J_\beta(A):=h_Z(A)-\beta\,\kappa(A)
\qquad (\beta\ge 0,\ A\ge 1).
\]
则连续 Pareto 前沿恰为
\[
A\in\left[1,\frac32\right].
\]
并且黄金端点 \(A=1\) 是一个暴露极点：存在且仅存在一条支持斜率
\[
\beta_c=
\frac{h_Z'(1)}{\kappa'(1)}
=
\frac{\left(\frac12-\frac1{\sqrt5}\right)(\log\varphi)^2}{\log\varphi-\frac1{\sqrt5}}
\approx 0.35953,
\]
使得
\[
A_\beta=1
\iff
\beta\ge \beta_c.
\]
当 \(0\le \beta<\beta_c\) 时，唯一极大点严格落在内部
\[
A_\beta\in\left(1,\frac32\right].
\]
\end{theorem}
\begin{proof}
命题 \ref{prop:metallic-pareto-frontier-lambda} 已证明：在“最大化 \(h_Z\)、最小化 \(\kappa\)”的双目标意义下，所有 \(A>3/2\) 的分支都被 \(A_\star=3/2\) 严格支配，而区间 \([1,3/2]\) 上任意两点互不支配，故连续 Pareto 前沿恰为 \([1,3/2]\)。

另一方面，命题 \ref{prop:metallic-linear-scalarization-threshold} 已给出端点阈值
\[
\beta_c=\frac{h_Z'(1)}{\kappa'(1)},
\]
并严格刻画了标量化极值点：当 \(\beta\ge \beta_c\) 时，\(A_\beta=1\)；当 \(0\le \beta<\beta_c\) 时，存在唯一内点解 \(A_\beta\in(1,3/2]\)。这正是黄金端点作为暴露极点的定义。证毕。
\end{proof}

\begin{corollary}[黄金端点的精确支撑超平面]\label{cor:conclusion-metallic-golden-supporting-hyperplane}
对一切 \(A\ge 1\)，成立全局仿射支撑不等式
\[
h_Z(A)-h_Z(1)\le
\beta_c\bigl(\kappa(A)-\kappa(1)\bigr),
\]
其中 \(\beta_c\) 由定理 \ref{thm:conclusion-metallic-pareto-exposed-golden-endpoint} 给出。并且等号当且仅当 \(A=1\)。
\end{corollary}
\begin{proof}
当 \(\beta=\beta_c\) 时，由定理 \ref{thm:conclusion-metallic-pareto-exposed-golden-endpoint}，端点 \(A=1\) 是 \(J_{\beta_c}(A)=h_Z(A)-\beta_c\kappa(A)\) 的唯一极大点。因此对任意 \(A\ge 1\)，
\[
h_Z(A)-\beta_c\kappa(A)\le h_Z(1)-\beta_c\kappa(1).
\]
移项即得所述不等式。唯一极大性迫使等号仅在 \(A=1\) 处成立。证毕。
\end{proof}

\begin{theorem}[有限样本连续标量化的离黄金端点律]\label{thm:conclusion-metallic-finite-sample-antigolden}
令
\[
\beta(N):=\frac{\log N}{N},
\qquad
N\in\ZZ_{\ge 1},
\]
并考虑连续标量化问题
\[
\max_{A\ge 1}J_{\beta(N)}(A).
\]
则对每个整数 \(N\ge 4\)，其唯一最优解满足
\[
A_{\beta(N)}\in\left(1,\frac32\right].
\]
特别地，黄金端点 \(A=1\) 在任何非平凡有限样本尺度 \(N\ge 4\) 上都不是连续最优解；并且
\[
A_{\beta(N)}\longrightarrow \frac32
\qquad (N\to\infty).
\]
\end{theorem}
\begin{proof}
由
\[
\beta(4)=\frac{\log 4}{4}\approx 0.34657<\beta_c
\]
以及 \(N\mapsto (\log N)/N\) 在 \(N\ge e\) 上严格递减，可知对一切 \(N\ge 4\) 都有
\[
\beta(N)<\beta_c.
\]
于是由定理 \ref{thm:conclusion-metallic-pareto-exposed-golden-endpoint}，最大化点必严格落在
\[
\left(1,\frac32\right].
\]
又因为 \(\beta(N)\downarrow 0\) 且命题 \ref{prop:metallic-linear-scalarization-threshold} 已给出 \(\beta\downarrow 0\) 时 \(A_\beta\to 3/2\)，故
\[
A_{\beta(N)}\to \frac32.
\]
证毕。
\end{proof}

\leanverified{paper\_conclusion\_metallic\_finite\_sample\_antigolden}

\begin{theorem}[整数设计空间的二值塌缩]\label{thm:conclusion-metallic-integer-binary-collapse}
\leanverified{paper\_conclusion\_metallic\_integer\_binary\_collapse}
在整数约束
\[
A\in\ZZ_{\ge 1}
\]
下，压缩指数
\[
h_Z(A)=\log\frac{A+1}{\lambda_A}
\]
的唯一最大值由
\[
A=2
\]
取得。并且对一切整数 \(A\ge 3\)，均被 \(A=2\) 在二目标 \((h_Z,-\kappa)\) 上严格支配。故整数 Pareto 集恰为
\[
\{1,2\}.
\]
\end{theorem}
\begin{proof}
命题 \ref{prop:metallic-compression-extremum} 已证明连续压缩极大点位于 \(A=3/2\)。因此整数点上的最大值只能落在相邻候选 \(A=1,2\) 之间。正文关于整数设计空间的缩减进一步给出：所有整数 \(A\ge 3\) 都被 \(A=2\) 在“最大化 \(h_Z\)、最小化 \(\kappa\)”的双目标意义下严格支配；这正是定理 \ref{thm:metallic-integer-scalarization-threshold} 的前提结构。于是整数 Pareto 集恰为 \(\{1,2\}\)。

再由
\[
h_Z(2)>h_Z(1)
\]
可知压缩指数的整数唯一极大点为 \(A=2\)。证毕。
\end{proof}


\begin{conjecture}[离散 kink normal fan 向连续 Legendre front 的收敛]\label{conj:conclusion-discrete-kink-normal-fan-hausdorff-to-continuous-legendre}
设附录 \ref{app:pom-projective-pressure-operator} 中的 Galerkin 或配点近似给出一列凸函数
\[
\Lambda_N:[0,\infty)\to\RR,
\]
满足：
\[
\Lambda_N\to\Lambda\quad\text{在每个紧区间上一致收敛};
\]
\[
\Lambda_N(q-1)=\log r_q-\log 2
\quad\text{对所有 }2\le q\le Q_N,\qquad Q_N\to\infty;
\]
并且每个 \(\Lambda_N\) 都保持凸性与整数锚点一致性。定义离散前沿
\[
G_N(c):=\min_{2\le q\le Q_N}\frac{\Lambda_N(q-1)+c}{q-1},
\]
以及连续前沿
\[
G(c):=\inf_{t>0}\frac{\Lambda(t)+c}{t}.
\]
则应有
\[
G_N\to G
\qquad\text{在每个有界 \(c\)-区间上一致收敛},
\]
并且 \(G_N\) 的 normal fan 在 Hausdorff 意义下收敛到 \(G\) 的切线扇。等价地，离散 kink 扇分相只是连续 Legendre 前沿的有限维折线近似；工程上“选 \(q\)”的动作最终应被重新解释为连续压力曲线切线法的整数采样。
\end{conjecture}


\paragraph{统一读法}
这一组结果把异常类、预算选择、连续压力接口、prime-register 证书前沿与 residue 细分压成了同一条精确链：异常的本体首先是一个商 Hilbert 空间，其大小由纯交互分量的正交能量唯一决定；两体审计只在“至多两个非平凡轴”这一退化情形下才完备，而一旦进入一般多轴情形，其盲区具有显式余维；预算最优化在离散层面上由 kink 集完全分相，但一旦接上附录 \ref{app:pom-projective-pressure-operator} 的连续压力曲线，便可看出离散 \(q\)-相图只是连续 Legendre 前沿的折线外壳。真正的连续相只活在整数节点的次微分上，而更大的 secant 相带内部还可能出现系统性的幽灵扇区；与此同时，证书前沿本身也不再应理解为“选一个最优整数矩阶”，而应理解为离散率函数的广义逆，以及其连续凸化闭包所给出的 Euler--Fenchel 切触曲线。在固定置信度口径下，连续最优预算的主指数趋向典型纤维指数，而首个非平凡修正则由 \(\Lambda''(0)\) 给出的中心极限型方差精确控制。到了 residue 轴上，纯度与均匀化给出 Jensen 压降带的两端刚性，而整个离散 Legendre 前沿又被严格夹在“原前沿”与“按 \((\log P)/m\) 平移后的原前沿”之间；二阶缺口与全部高阶细分矩则进一步塌缩为局部 Fourier 模的精确账本。这说明复杂性在这里并不首先体现在“需要更高阶矩”，而首先体现在局部交互类如何进入商空间几何、离散相图如何作为连续压力的 secant 包出现、prime-register 前沿如何作为连续 Legendre 曲线的整数采样，以及一阶局部层析如何生成整条高阶矩层级。在金属平均族这一一参数模型上，同一条 Legendre/Pareto 几何又进一步暴露出黄金端点的精确支撑超平面、有限样本惩罚下连续最优点对黄金端点的系统性离开，以及整数约束下向 \(\{1,2\}\) 的二值塌缩。

\endinput
